\documentclass[12pt,a4paper]{article}

\usepackage{fontspec}
\usepackage[brazilian]{babel}
\setmainfont{Arial}
\setsansfont{Arial}
\setmonofont{Consolas}
\renewcommand{\familydefault}{\sfdefault}

\usepackage{
  geometry,
  setspace,
  longtable,
  booktabs,
  array,
  calc,
  etoolbox,
  graphicx,
  enumitem,
  ragged2e,
  fancyvrb,
  fvextra,
  framed
}
\usepackage[table]{xcolor}
\usepackage{titlesec,fancyhdr,xurl,hyperref,bookmark}

\definecolor{AzulCorporativo}{RGB}{0,82,155}
\definecolor{AzulPetroleo}{RGB}{23,104,133}
\definecolor{CinzaEscuro}{RGB}{64,64,64}
\definecolor{CinzaClaro}{RGB}{242,242,242}
\definecolor{LaranjaDestaque}{RGB}{230,115,0}

\hypersetup{
  colorlinks=true,
  linkcolor=AzulPetroleo,
  urlcolor=AzulCorporativo,
  citecolor=LaranjaDestaque,
  pdftitle={Camada Analítica e Extração de Inteligência},
  pdfauthor={Lucas Dias Noronha}
}

\titleformat{\section}
  {\color{AzulCorporativo}\normalfont\Large\bfseries}
  {\thesection}{1em}{}
\titleformat{\subsection}
  {\color{AzulPetroleo}\normalfont\large\bfseries}
  {\thesubsection}{1em}{}
\titleformat{\subsubsection}
  {\color{CinzaEscuro}\normalfont\normalsize\bfseries}
  {\thesubsubsection}{1em}{}

\geometry{left=3cm,right=2cm,top=3cm,bottom=2cm,headheight=45pt}
\onehalfspacing
\setlength{\parskip}{0.35em}
\setlength{\emergencystretch}{3em}
\setlength{\LTleft}{0pt}
\setlength{\LTright}{0pt}
\setcounter{tocdepth}{2}
\newcounter{none}
\AtBeginEnvironment{longtable}{\footnotesize\setstretch{1.0}}

\pagestyle{fancy}
\fancyhf{}
\renewcommand{\headrulewidth}{0pt}
\fancyhead[C]{
  \colorbox{AzulCorporativo}{
    \makebox[\dimexpr\textwidth-2\fboxsep\relax][l]{
      \textcolor{white}{\textbf{\ \ Camada Analítica e Extração de Inteligência}}
    }
  }
}
\fancyfoot[C]{\color{CinzaEscuro}\thepage}

\providecommand{\tightlist}{
  \setlength{\itemsep}{0pt}\setlength{\parskip}{0pt}
}
\newcommand{\passthrough}[1]{#1}
\usepackage{color}
\usepackage{fancyvrb}
\newcommand{\VerbBar}{|}
\newcommand{\VERB}{\Verb[commandchars=\\\{\}]}
\DefineVerbatimEnvironment{Highlighting}{Verbatim}{commandchars=\\\{\}}
% Add ',fontsize=\small' for more characters per line
\newenvironment{Shaded}{}{}
\newcommand{\AlertTok}[1]{\textcolor[rgb]{1.00,0.00,0.00}{\textbf{#1}}}
\newcommand{\AnnotationTok}[1]{\textcolor[rgb]{0.38,0.63,0.69}{\textbf{\textit{#1}}}}
\newcommand{\AttributeTok}[1]{\textcolor[rgb]{0.49,0.56,0.16}{#1}}
\newcommand{\BaseNTok}[1]{\textcolor[rgb]{0.25,0.63,0.44}{#1}}
\newcommand{\BuiltInTok}[1]{\textcolor[rgb]{0.00,0.50,0.00}{#1}}
\newcommand{\CharTok}[1]{\textcolor[rgb]{0.25,0.44,0.63}{#1}}
\newcommand{\CommentTok}[1]{\textcolor[rgb]{0.38,0.63,0.69}{\textit{#1}}}
\newcommand{\CommentVarTok}[1]{\textcolor[rgb]{0.38,0.63,0.69}{\textbf{\textit{#1}}}}
\newcommand{\ConstantTok}[1]{\textcolor[rgb]{0.53,0.00,0.00}{#1}}
\newcommand{\ControlFlowTok}[1]{\textcolor[rgb]{0.00,0.44,0.13}{\textbf{#1}}}
\newcommand{\DataTypeTok}[1]{\textcolor[rgb]{0.56,0.13,0.00}{#1}}
\newcommand{\DecValTok}[1]{\textcolor[rgb]{0.25,0.63,0.44}{#1}}
\newcommand{\DocumentationTok}[1]{\textcolor[rgb]{0.73,0.13,0.13}{\textit{#1}}}
\newcommand{\ErrorTok}[1]{\textcolor[rgb]{1.00,0.00,0.00}{\textbf{#1}}}
\newcommand{\ExtensionTok}[1]{#1}
\newcommand{\FloatTok}[1]{\textcolor[rgb]{0.25,0.63,0.44}{#1}}
\newcommand{\FunctionTok}[1]{\textcolor[rgb]{0.02,0.16,0.49}{#1}}
\newcommand{\ImportTok}[1]{\textcolor[rgb]{0.00,0.50,0.00}{\textbf{#1}}}
\newcommand{\InformationTok}[1]{\textcolor[rgb]{0.38,0.63,0.69}{\textbf{\textit{#1}}}}
\newcommand{\KeywordTok}[1]{\textcolor[rgb]{0.00,0.44,0.13}{\textbf{#1}}}
\newcommand{\NormalTok}[1]{#1}
\newcommand{\OperatorTok}[1]{\textcolor[rgb]{0.40,0.40,0.40}{#1}}
\newcommand{\OtherTok}[1]{\textcolor[rgb]{0.00,0.44,0.13}{#1}}
\newcommand{\PreprocessorTok}[1]{\textcolor[rgb]{0.74,0.48,0.00}{#1}}
\newcommand{\RegionMarkerTok}[1]{#1}
\newcommand{\SpecialCharTok}[1]{\textcolor[rgb]{0.25,0.44,0.63}{#1}}
\newcommand{\SpecialStringTok}[1]{\textcolor[rgb]{0.73,0.40,0.53}{#1}}
\newcommand{\StringTok}[1]{\textcolor[rgb]{0.25,0.44,0.63}{#1}}
\newcommand{\VariableTok}[1]{\textcolor[rgb]{0.10,0.09,0.49}{#1}}
\newcommand{\VerbatimStringTok}[1]{\textcolor[rgb]{0.25,0.44,0.63}{#1}}
\newcommand{\WarningTok}[1]{\textcolor[rgb]{0.38,0.63,0.69}{\textbf{\textit{#1}}}}
\RecustomVerbatimEnvironment{Highlighting}{Verbatim}{
  commandchars=\\\{\},breaklines,breakanywhere,fontsize=\small
}
\DefineVerbatimEnvironment{verbatim}{Verbatim}{
  breaklines,breakanywhere,fontsize=\small
}

\begin{document}

\thispagestyle{empty}
\vspace*{\fill}
\begin{center}
  {\Huge\textcolor{AzulCorporativo}{\textbf{Camada Analítica e Extração
de Inteligência}}}\\[0.5cm]
  {\Large\textcolor{AzulPetroleo}{Modelo Analítico de Dados para E-commerce}}\\[0.8cm]
\end{center}
\vspace*{\fill}
\noindent
\textbf{\textcolor{AzulCorporativo}{Autor:}} Lucas Dias Noronha
\hfill
\textbf{\textcolor{AzulCorporativo}{Data:}} 2 de setembro de 2026

\newpage
\tableofcontents
\newpage

\section{Documentação
analítica}\label{docs__analytics__readme.md__documentauxe7uxe3o-analuxedtica}

Documentação destinada à leitura humana sobre métricas, decisões
analíticas, premissas, limitações e matrizes de rastreabilidade.

\subsection{Artefatos}\label{docs__analytics__readme.md__artefatos}

\begin{itemize}
\item
  \href{camada_analitica.tex}{\nolinkurl{camada_analitica.tex}} e
  \href{camada_analitica.pdf}{\nolinkurl{camada_analitica.pdf}}:
  documentação consolidada da camada analítica, reunindo os artefatos
  Markdown abaixo em um único documento.
\item
  \href{contrato_analitico.md}{\nolinkurl{contrato_analitico.md}}:
  problema de negócio, stakeholders, demandas, casos de uso, catálogo
  preliminar de indicadores, regras semânticas e limitações da camada
  analítica.
\item
  \href{evidencias_consultas_basicas_core.md}{\nolinkurl{evidencias_consultas_basicas_core.md}}:
  resultados de controle e decisões de granularidade observados durante
  a etapa M06-01.
\item
  \href{evidencias_consultas_analiticas_core.md}{\nolinkurl{evidencias_consultas_analiticas_core.md}}:
  resultados financeiros, concentração de vendedores, pagamentos e
  riscos associados observados durante a etapa M06-02.
\item
  \href{matriz_validacao_requisitos_funcionais.md}{\nolinkurl{matriz_validacao_requisitos_funcionais.md}}:
  rastreabilidade entre RF01--RF12, consultas e evidências da etapa
  M06-03.
\item
  \href{catalogo_metricas_analiticas.md}{\nolinkurl{catalogo_metricas_analiticas.md}}:
  definições, fórmulas, granularidades, resultados e limitações das
  métricas financeiras e de vendedores consolidadas na etapa M06-04.
\item
  \href{estruturas_schema_analytics.md}{\nolinkurl{estruturas_schema_analytics.md}}:
  finalidade, granularidade, campos, dependências e validação das views
  criadas na M06-05.
\item
  \href{consumo_power_bi_sql.md}{\nolinkurl{consumo_power_bi_sql.md}}:
  contrato de datasets, relacionamentos, medidas, atualização e
  controles para consumo no Power BI e por SQL, consolidado na M06-06.
\end{itemize}

\section{Contrato analítico da camada
ANALYTICS}\label{docs__analytics__contrato_analitico.md__contrato-analuxedtico-da-camada-analytics}

\subsection{Finalidade}\label{docs__analytics__contrato_analitico.md__finalidade}

Este documento define o problema de negócio, os stakeholders, as
decisões, as demandas, os casos de uso e o catálogo preliminar de
indicadores que orientam a fase M06 --- Camada Analítica e Extração de
Inteligência.

O contrato antecede a implementação das consultas e estruturas
analíticas. Ele deve orientar a exploração da CORE sem transformar as
issues da fase em uma sequência mecânica. Exploração, validação de
requisitos e definição de métricas podem se retroalimentar, desde que
alterações sejam sustentadas por evidência e registradas de forma
rastreável.

\subsection{Natureza do
cenário}\label{docs__analytics__contrato_analitico.md__natureza-do-cenuxe1rio}

O Brazilian E-Commerce Public Dataset by Olist não identifica uma
empresa demandante nem sua estrutura organizacional. Por isso, os
stakeholders e as decisões descritos neste documento constituem
\textbf{personas analíticas hipotéticas}, utilizadas para organizar um
caso de uso coerente com os dados disponíveis.

Os resultados não devem ser apresentados como diagnóstico da operação
atual da Olist nem como representação de sua estrutura real de gestão.

\subsection{Cenário e problema de
negócio}\label{docs__analytics__contrato_analitico.md__cenuxe1rio-e-problema-de-neguxf3cio}

Uma empresa operadora de marketplace precisa compreender seu desempenho
histórico para orientar decisões financeiras, comerciais e operacionais.
Embora possua dados integrados sobre pedidos, clientes, itens, produtos,
vendedores, pagamentos, avaliações e entregas, ainda não dispõe de uma
visão analítica consolidada e rastreável.

O problema de negócio adotado é:

\begin{quote}
A gestão do marketplace não dispõe de informações consolidadas e
rastreáveis para avaliar prioritariamente o desempenho financeiro e a
contribuição dos vendedores, considerando também os riscos logísticos e
de experiência do cliente associados aos resultados observados.
\end{quote}

O produto analítico apoiará a compreensão e a priorização de decisões. O
dataset não oferece informações suficientes para executar ações
operacionais ou atribuir causalidade aos padrões encontrados.

\subsection{Objetivo
analítico}\label{docs__analytics__contrato_analitico.md__objetivo-analuxedtico}

Construir uma camada analítica que transforme os dados históricos do
marketplace em indicadores financeiros e de desempenho dos vendedores,
permitindo:

\begin{itemize}
\tightlist
\item
  acompanhar a geração, a composição e a evolução dos valores
  observados;
\item
  reconciliar as representações financeiras presentes em itens e
  pagamentos;
\item
  identificar vendedores relevantes e mensurar a concentração da base;
\item
  comparar categorias, regiões e períodos;
\item
  contextualizar o desempenho com aspectos logísticos e de experiência;
\item
  preparar estruturas reproduzíveis para consumo por SQL e ferramentas
  de BI.
\end{itemize}

A pergunta executiva central é:

\begin{quote}
Quanto valor foi movimentado em pedidos entregues, como esse valor
evoluiu e foi composto, quais vendedores mais contribuíram e quais
riscos operacionais ou de experiência estão associados a essa
contribuição?
\end{quote}

\subsection{Pilares do produto
analítico}\label{docs__analytics__contrato_analitico.md__pilares-do-produto-analuxedtico}

\subsubsection{Desempenho financeiro --- foco
principal}\label{docs__analytics__contrato_analitico.md__desempenho-financeiro-foco-principal}

Compreender a geração, a composição, a evolução e a concentração dos
valores observados no marketplace. Inclui valores dos itens, frete,
pagamentos, parcelamento, reconciliação e recortes por período,
categoria, vendedor e região.

\subsubsection{Desempenho dos vendedores --- foco
secundário}\label{docs__analytics__contrato_analitico.md__desempenho-dos-vendedores-foco-secunduxe1rio}

Avaliar a contribuição financeira, a concentração, o alcance comercial e
a qualidade operacional associada aos vendedores.

\subsubsection{Logística e experiência --- pilares
complementares}\label{docs__analytics__contrato_analitico.md__loguxedstica-e-experiuxeancia-pilares-complementares}

Utilizar prazos, atrasos, frete e avaliações para contextualizar riscos
e diferenças de desempenho. Associação estatística ou descritiva não
será tratada como evidência de causalidade.

\subsection{Stakeholders e decisões
apoiadas}\label{docs__analytics__contrato_analitico.md__stakeholders-e-decisuxf5es-apoiadas}

{\def\LTcaptype{none} % do not increment counter
\begin{longtable}[]{@{}
  >{\raggedleft\arraybackslash}p{(\linewidth - 6\tabcolsep) * \real{0.3077}}
  >{\raggedright\arraybackslash}p{(\linewidth - 6\tabcolsep) * \real{0.2308}}
  >{\raggedright\arraybackslash}p{(\linewidth - 6\tabcolsep) * \real{0.2308}}
  >{\raggedright\arraybackslash}p{(\linewidth - 6\tabcolsep) * \real{0.2308}}@{}}
\toprule\noalign{}
\begin{minipage}[b]{\linewidth}\raggedleft
Prioridade
\end{minipage} & \begin{minipage}[b]{\linewidth}\raggedright
Stakeholder hipotético
\end{minipage} & \begin{minipage}[b]{\linewidth}\raggedright
Interesse
\end{minipage} & \begin{minipage}[b]{\linewidth}\raggedright
Decisões apoiadas
\end{minipage} \\
\midrule\noalign{}
\endhead
\bottomrule\noalign{}
\endlastfoot
1 & Gestão financeira e executiva & Valores, composição, evolução e
concentração & Priorizar períodos, segmentos e riscos para
investigação \\
2 & Gestão de vendedores & Contribuição e qualidade da base & Priorizar
acompanhamento e desenvolvimento de vendedores \\
3 & Gestão comercial & Categorias, regiões e comportamento dos pedidos &
Identificar fontes de valor e segmentos relevantes \\
4 & Operações e logística & Frete, prazos e atrasos & Priorizar
investigações operacionais \\
5 & Experiência do cliente & Avaliações e fatores associados &
Identificar jornadas e segmentos insatisfatórios \\
6 & BI e dados & Definições, rastreabilidade e consumo & Disponibilizar
indicadores reproduzíveis \\
\end{longtable}
}

\subsection{Demandas e perguntas
analíticas}\label{docs__analytics__contrato_analitico.md__demandas-e-perguntas-analuxedticas}

{\def\LTcaptype{none} % do not increment counter
\begin{longtable}[]{@{}
  >{\raggedleft\arraybackslash}p{(\linewidth - 8\tabcolsep) * \real{0.2500}}
  >{\raggedright\arraybackslash}p{(\linewidth - 8\tabcolsep) * \real{0.1875}}
  >{\raggedright\arraybackslash}p{(\linewidth - 8\tabcolsep) * \real{0.1875}}
  >{\raggedright\arraybackslash}p{(\linewidth - 8\tabcolsep) * \real{0.1875}}
  >{\raggedright\arraybackslash}p{(\linewidth - 8\tabcolsep) * \real{0.1875}}@{}}
\toprule\noalign{}
\begin{minipage}[b]{\linewidth}\raggedleft
Prioridade
\end{minipage} & \begin{minipage}[b]{\linewidth}\raggedright
Stakeholder
\end{minipage} & \begin{minipage}[b]{\linewidth}\raggedright
Decisão
\end{minipage} & \begin{minipage}[b]{\linewidth}\raggedright
Pergunta
\end{minipage} & \begin{minipage}[b]{\linewidth}\raggedright
Indicadores principais
\end{minipage} \\
\midrule\noalign{}
\endhead
\bottomrule\noalign{}
\endlastfoot
P0 & Financeiro/executivo & Acompanhar desempenho & Quanto foi
movimentado e como evoluiu? & Valores dos itens, frete, valor bruto,
valor pago e crescimento \\
P0 & Financeiro/executivo & Validar consistência & O valor pago é
compatível com itens mais frete? & Diferenças absoluta e percentual \\
P0 & Gestão de vendedores & Identificar vendedores estratégicos & Quais
vendedores mais contribuem? & Valor dos itens, participação e posição \\
P0 & Gestão de vendedores & Avaliar dependência & O resultado está
concentrado? & Top N, curva ABC e HHI \\
P1 & Financeiro & Entender a composição & Qual é a participação do
frete? & Frete e participação no valor bruto \\
P1 & Financeiro & Entender pagamentos & Quais meios e condições são
utilizados? & Valor e pedidos por tipo e parcelamento \\
P1 & Comercial & Identificar fontes de valor & Quais categorias e
regiões contribuem mais? & Valor, pedidos, itens e participação \\
P1 & Gestão de vendedores & Avaliar produtividade & Quem combina valor,
pedidos e itens? & Valores, volumes e médias por vendedor \\
P1 & Gestão de vendedores & Avaliar qualidade associada & Quem participa
de pedidos com atraso ou avaliação ruim? & Taxas associadas de atraso e
avaliação negativa \\
P2 & Operações & Investigar riscos logísticos & Onde estão atrasos e
fretes elevados? & Taxa de atraso, prazo e frete \\
P2 & Experiência & Investigar insatisfação & Como avaliações variam por
entrega e segmento? & Nota média, distribuição e avaliações negativas \\
P2 & BI e dados & Preparar consumo recorrente & Quais definições
precisam ser estabilizadas? & Catálogo, queries, views e marts \\
\end{longtable}
}

\subsection{Casos de uso
priorizados}\label{docs__analytics__contrato_analitico.md__casos-de-uso-priorizados}

\subsubsection{Produto
mínimo}\label{docs__analytics__contrato_analitico.md__produto-muxednimo}

\paragraph{UC01 --- Resumo financeiro
executivo}\label{docs__analytics__contrato_analitico.md__uc01-resumo-financeiro-executivo}

Apresentar, por mês, pedidos entregues, valor dos itens, frete, valor
bruto, valor pago registrado, diferença de reconciliação, valor médio
por pedido e crescimento.

\paragraph{UC02 --- Composição dos
pagamentos}\label{docs__analytics__contrato_analitico.md__uc02-composiuxe7uxe3o-dos-pagamentos}

Analisar participação por tipo de pagamento, valor médio por modalidade,
parcelamento, pedidos com múltiplas formas de pagamento e ocorrências
não definidas.

\paragraph{UC03 --- Desempenho financeiro dos
vendedores}\label{docs__analytics__contrato_analitico.md__uc03-desempenho-financeiro-dos-vendedores}

Analisar vendedores ativos, valores, itens, pedidos, participação
individual e acumulada, rankings, curva ABC e concentração.

\paragraph{UC04 --- Reconciliação
financeira}\label{docs__analytics__contrato_analitico.md__uc04-reconciliauxe7uxe3o-financeira}

Comparar o valor bruto calculado nos itens com o valor registrado nos
pagamentos, preservando a granularidade de pedido e identificando
exceções.

\subsubsection{Diagnósticos
importantes}\label{docs__analytics__contrato_analitico.md__diagnuxf3sticos-importantes}

\paragraph{UC05 --- Perfil e alcance dos
vendedores}\label{docs__analytics__contrato_analitico.md__uc05-perfil-e-alcance-dos-vendedores}

Analisar categorias comercializadas, origem, regiões atendidas,
amplitude geográfica e dependência de categorias ou regiões.

\paragraph{UC06 --- Risco operacional associado aos
vendedores}\label{docs__analytics__contrato_analitico.md__uc06-risco-operacional-associado-aos-vendedores}

Combinar contribuição financeira, envio após o limite, atraso, frete e
avaliação. O resultado deve separar responsabilidade direta de mera
associação.

\paragraph{UC07 --- Categorias e regiões
financeiras}\label{docs__analytics__contrato_analitico.md__uc07-categorias-e-regiuxf5es-financeiras}

Identificar segmentos com maior valor, evolução mensal, valor médio,
peso do frete, concentração de vendedores e qualidade logística
associada.

\subsubsection{Análises condicionadas à
evidência}\label{docs__analytics__contrato_analitico.md__anuxe1lises-condicionadas-uxe0-eviduxeancia}

\begin{itemize}
\tightlist
\item
  segmentação de vendedores em quadrantes;
\item
  curva ABC por período;
\item
  HHI mensal;
\item
  estimativas geográficas de distância;
\item
  estruturas persistentes específicas para BI.
\end{itemize}

Essas análises somente serão consolidadas se a distribuição dos dados e
a necessidade de consumo justificarem sua implementação.

\subsection{Política
financeira}\label{docs__analytics__contrato_analitico.md__poluxedtica-financeira}

\subsubsection{Data de
referência}\label{docs__analytics__contrato_analitico.md__data-de-referuxeancia}

A dimensão temporal principal das métricas financeiras e comerciais será
\nolinkurl{core.pedido.data_compra}. As demais datas serão utilizadas
para finalidades específicas:

\begin{itemize}
\tightlist
\item
  \nolinkurl{data_aprovacao}: processo de aprovação;
\item
  \nolinkurl{data_envio_transportador}: processo logístico;
\item
  \nolinkurl{data_entrega}: realização da entrega;
\item
  \nolinkurl{data_estimada}: cumprimento do prazo.
\end{itemize}

Uma métrica não deve alternar entre datas sem registrar explicitamente a
mudança de significado.

\subsubsection{Política de
status}\label{docs__analytics__contrato_analitico.md__poluxedtica-de-status}

A visão principal utilizará pedidos com
\nolinkurl{status_pedido = 'delivered'}. Essa base representa pedidos
concluídos, mas não comprova reconhecimento contábil de receita.

Uma visão ampla poderá incluir todos os pedidos com registros
financeiros para reconciliação e investigação. Pedidos cancelados,
indisponíveis ou ainda não concluídos serão analisados separadamente.

\subsubsection{Medidas
observadas}\label{docs__analytics__contrato_analitico.md__medidas-observadas}

{\def\LTcaptype{none} % do not increment counter
\begin{longtable}[]{@{}
  >{\raggedright\arraybackslash}p{(\linewidth - 8\tabcolsep) * \real{0.2000}}
  >{\raggedright\arraybackslash}p{(\linewidth - 8\tabcolsep) * \real{0.2000}}
  >{\raggedright\arraybackslash}p{(\linewidth - 8\tabcolsep) * \real{0.2000}}
  >{\raggedright\arraybackslash}p{(\linewidth - 8\tabcolsep) * \real{0.2000}}
  >{\raggedright\arraybackslash}p{(\linewidth - 8\tabcolsep) * \real{0.2000}}@{}}
\toprule\noalign{}
\begin{minipage}[b]{\linewidth}\raggedright
Métrica
\end{minipage} & \begin{minipage}[b]{\linewidth}\raggedright
Fórmula
\end{minipage} & \begin{minipage}[b]{\linewidth}\raggedright
Unidade
\end{minipage} & \begin{minipage}[b]{\linewidth}\raggedright
Granularidade de origem
\end{minipage} & \begin{minipage}[b]{\linewidth}\raggedright
Fonte
\end{minipage} \\
\midrule\noalign{}
\endhead
\bottomrule\noalign{}
\endlastfoot
Valor dos itens & \nolinkurl{SUM(preco_item)} & BRL & Item &
\nolinkurl{core.item_pedido.preco_item} \\
Valor de frete & \nolinkurl{SUM(valor_frete)} & BRL & Item &
\nolinkurl{core.item_pedido.valor_frete} \\
Valor bruto & \nolinkurl{SUM(preco_item + valor_frete)} & BRL & Item,
agregado por pedido & \nolinkurl{core.item_pedido} \\
Valor pago registrado & \nolinkurl{SUM(valor_pagamento)} & BRL &
Sequência, agregada por pedido &
\nolinkurl{core.pagamento.valor_pagamento} \\
Diferença de reconciliação & \nolinkurl{valor_pago - valor_bruto} & BRL
& Pedido & Itens e pagamentos pré-agregados \\
Diferença percentual &
\nolinkurl{(valor_pago - valor_bruto) / valor_bruto} & Percentual &
Pedido & Itens e pagamentos pré-agregados \\
Participação do frete & \nolinkurl{valor_frete / valor_bruto} &
Percentual & Período ou segmento & \nolinkurl{core.item_pedido} \\
Valor médio por pedido & \nolinkurl{valor agregado / pedidos distintos}
& BRL/pedido & Período ou segmento & Pedido e medida escolhida \\
Itens por pedido & \nolinkurl{itens / pedidos distintos} & Itens/pedido
& Período ou segmento & \nolinkurl{core.item_pedido} \\
\end{longtable}
}

A medida utilizada no numerador do valor médio deve aparecer em seu nome
ou em sua definição. Não será adotado um ``ticket médio'' sem
especificar se a base é o valor dos itens, o valor bruto ou o valor
pago.

\subsubsection{Reconciliação}\label{docs__analytics__contrato_analitico.md__reconciliauxe7uxe3o}

\nolinkurl{core.item_pedido} e \nolinkurl{core.pagamento} devem ser
agregadas separadamente por \nolinkurl{id_pedido} antes de qualquer
combinação. A existência de divergências não autoriza concluir que uma
representação esteja incorreta sem investigação.

\subsubsection{Atribuição aos
vendedores}\label{docs__analytics__contrato_analitico.md__atribuiuxe7uxe3o-aos-vendedores}

Podem ser atribuídos diretamente ao vendedor:

\begin{itemize}
\tightlist
\item
  valor dos itens;
\item
  frete associado aos itens;
\item
  valor bruto dos itens;
\item
  quantidade de itens;
\item
  quantidade de pedidos distintos com participação.
\end{itemize}

O valor de pagamento pertence ao pedido. Ele não será atribuído
integralmente a cada vendedor. Um eventual rateio proporcional deverá
ser nomeado como \textbf{valor alocado}, acompanhado da fórmula e sem
ser confundido com pagamento observado.

\subsection{Política de análise dos
vendedores}\label{docs__analytics__contrato_analitico.md__poluxedtica-de-anuxe1lise-dos-vendedores}

\subsubsection{Métricas
fundamentais}\label{docs__analytics__contrato_analitico.md__muxe9tricas-fundamentais}

{\def\LTcaptype{none} % do not increment counter
\begin{longtable}[]{@{}
  >{\raggedright\arraybackslash}p{(\linewidth - 4\tabcolsep) * \real{0.3333}}
  >{\raggedright\arraybackslash}p{(\linewidth - 4\tabcolsep) * \real{0.3333}}
  >{\raggedright\arraybackslash}p{(\linewidth - 4\tabcolsep) * \real{0.3333}}@{}}
\toprule\noalign{}
\begin{minipage}[b]{\linewidth}\raggedright
Métrica
\end{minipage} & \begin{minipage}[b]{\linewidth}\raggedright
Fórmula resumida
\end{minipage} & \begin{minipage}[b]{\linewidth}\raggedright
Interpretação
\end{minipage} \\
\midrule\noalign{}
\endhead
\bottomrule\noalign{}
\endlastfoot
Vendedores com pedidos & \nolinkurl{COUNT(DISTINCT id_vendedor)} &
Vendedores presentes na base selecionada \\
Vendedores ativos & Vendedores com item em pedido entregue no período &
Base efetivamente ativa segundo a definição adotada \\
Pedidos por vendedor & \nolinkurl{COUNT(DISTINCT id_pedido)} & Pedidos
com participação \\
Itens por vendedor & Contagem de itens & Volume comercializado \\
Valor dos itens & Soma de \nolinkurl{preco_item} & Contribuição
comercial direta \\
Frete associado & Soma de \nolinkurl{valor_frete} & Frete cobrado nos
itens \\
Valor bruto & Itens mais frete & Valor bruto associado \\
Valor médio por item & Valor dos itens dividido pelos itens & Perfil de
preço \\
Valor médio por pedido com participação & Valor dos itens dividido pelos
pedidos distintos & Contribuição média nos pedidos \\
Participação financeira & Valor do vendedor dividido pelo total &
Relevância relativa \\
Categorias comercializadas & Categorias distintas & Diversidade
observada \\
Estados atendidos & Estados distintos dos clientes & Alcance geográfico
observado \\
\end{longtable}
}

\subsubsection{Concentração}\label{docs__analytics__contrato_analitico.md__concentrauxe7uxe3o}

Devem ser avaliadas a participação individual e as participações
acumuladas dos 5, 10 e 20 maiores vendedores e dos 10\% maiores
vendedores.

A curva ABC poderá utilizar inicialmente os limites acumulados de 80\%
para a classe A e 95\% para as classes A+B. Esses limites são parâmetros
descritivos e podem ser revistos após a análise da distribuição.

O índice HHI será calculado pela soma dos quadrados das participações.
Ele mede concentração na base observada e não deve ser interpretado como
concentração do mercado brasileiro nem como avaliação regulatória.

\subsubsection{Indicadores
operacionais}\label{docs__analytics__contrato_analitico.md__indicadores-operacionais}

São diretamente atribuíveis ao vendedor:

\begin{itemize}
\tightlist
\item
  cumprimento da \nolinkurl{data_limite_envio} do item;
\item
  valor dos itens e frete;
\item
  categorias comercializadas;
\item
  origem e alcance geográfico observados.
\end{itemize}

São apenas associados ao vendedor:

\begin{itemize}
\tightlist
\item
  atraso da entrega final;
\item
  tempo total de entrega;
\item
  avaliação;
\item
  cancelamento do pedido.
\end{itemize}

Quando houver múltiplos vendedores no mesmo pedido, deve-se utilizar a
expressão ``pedido associado'' ou ``pedido com participação do
vendedor''.

\subsubsection{Segmentação
diagnóstica}\label{docs__analytics__contrato_analitico.md__segmentauxe7uxe3o-diagnuxf3stica}

Uma análise posterior poderá classificar vendedores nos quadrantes
estratégico, atenção prioritária, potencial e baixa contribuição
crítica, combinando desempenho financeiro e qualidade operacional. Os
limites de alto, médio e baixo serão definidos somente após examinar a
distribuição, por percentis ou faixas justificadas.

\subsection{Catálogo mínimo de
KPIs}\label{docs__analytics__contrato_analitico.md__catuxe1logo-muxednimo-de-kpis}

\subsubsection{Financeiros}\label{docs__analytics__contrato_analitico.md__financeiros}

\begin{enumerate}
\def\labelenumi{\arabic{enumi}.}
\tightlist
\item
  Valor dos itens de pedidos entregues.
\item
  Valor de frete de pedidos entregues.
\item
  Valor bruto de pedidos entregues.
\item
  Valor pago registrado em pedidos entregues.
\item
  Diferença de reconciliação financeira.
\item
  Valor médio por pedido entregue, com medida-base explícita.
\item
  Participação do frete no valor bruto.
\item
  Crescimento mensal do valor dos itens.
\item
  Taxa de cancelamento por quantidade.
\item
  Valor associado a pedidos cancelados.
\end{enumerate}

\subsubsection{Vendedores}\label{docs__analytics__contrato_analitico.md__vendedores}

\begin{enumerate}
\def\labelenumi{\arabic{enumi}.}
\setcounter{enumi}{10}
\tightlist
\item
  Vendedores ativos.
\item
  Valor dos itens por vendedor.
\item
  Participação financeira por vendedor.
\item
  Participação acumulada dos maiores vendedores.
\item
  Pedidos com participação do vendedor.
\item
  Itens por vendedor.
\item
  Valor médio do vendedor por pedido com participação.
\item
  Taxa de itens enviados após o limite.
\item
  Taxa de pedidos atrasados associados ao vendedor.
\item
  Nota média dos pedidos associados ao vendedor.
\end{enumerate}

O catálogo é preliminar. A exploração deve validar disponibilidade,
qualidade, distribuição e utilidade antes de promover uma métrica a
estrutura persistente ou KPI de consumo.

\subsection{Limitações}\label{docs__analytics__contrato_analitico.md__limitauxe7uxf5es}

\subsubsection{Financeiras}\label{docs__analytics__contrato_analitico.md__financeiras}

O dataset não contém custos dos produtos, comissões, impostos, descontos
explicitamente separados, tarifas de pagamento, custos logísticos reais,
estornos completos, margem ou lucro. Valores registrados não comprovam
liquidação contábil.

\subsubsection{Temporais}\label{docs__analytics__contrato_analitico.md__temporais}

Os dados representam um período histórico encerrado. Não permitem
monitoramento em tempo real e não representam a operação atual. Meses
incompletos devem ser identificados antes de comparações ou cálculos de
crescimento.

\subsubsection{Vendedores}\label{docs__analytics__contrato_analitico.md__vendedores-1}

Um pedido pode envolver vários vendedores. Pagamentos e avaliações
pertencem ao pedido. O atraso final também pode depender de fatores não
representados. Não há custos, contratos, comissões ou metas individuais.

\subsubsection{Clientes e
mercado}\label{docs__analytics__contrato_analitico.md__clientes-e-mercado}

Somente clientes com transações observadas estão presentes. Não há
visitantes, abandono, conversão ou dimensão completa do mercado.
Participação no dataset não equivale a participação de mercado.

\subsubsection{Logística e
geografia}\label{docs__analytics__contrato_analitico.md__loguxedstica-e-geografia}

Não há transportadora, rota, veículo nem rastreamento detalhado.
Coordenadas são associadas a prefixos de CEP e podem possuir múltiplos
registros. Distâncias derivadas serão aproximações geográficas, não
distâncias percorridas.

\subsection{Convenções de
nomenclatura}\label{docs__analytics__contrato_analitico.md__convenuxe7uxf5es-de-nomenclatura}

{\def\LTcaptype{none} % do not increment counter
\begin{longtable}[]{@{}
  >{\raggedright\arraybackslash}p{(\linewidth - 2\tabcolsep) * \real{0.5000}}
  >{\raggedright\arraybackslash}p{(\linewidth - 2\tabcolsep) * \real{0.5000}}@{}}
\toprule\noalign{}
\begin{minipage}[b]{\linewidth}\raggedright
Evitar
\end{minipage} & \begin{minipage}[b]{\linewidth}\raggedright
Utilizar
\end{minipage} \\
\midrule\noalign{}
\endhead
\bottomrule\noalign{}
\endlastfoot
Receita & Valor dos itens ou valor movimentado, com definição \\
Faturamento & Valor bruto dos pedidos \\
Lucro & Não disponível \\
Rentabilidade do vendedor & Desempenho financeiro do vendedor \\
Ticket médio do vendedor & Valor médio do vendedor por pedido com
participação \\
Custo de frete & Valor de frete cobrado \\
Vendas concluídas & Pedidos entregues \\
Atraso do vendedor & Pedido atrasado associado ao vendedor \\
Cliente ativo & Cliente com pedido no período, quando definido \\
Participação de mercado & Participação no valor observado no dataset \\
\end{longtable}
}

\subsection{Regras de qualidade
analítica}\label{docs__analytics__contrato_analitico.md__regras-de-qualidade-analuxedtica}

Toda métrica consolidada deve informar, quando aplicável:

\begin{itemize}
\tightlist
\item
  nome e objetivo;
\item
  fórmula e unidade;
\item
  granularidade;
\item
  data de referência;
\item
  filtros e status incluídos;
\item
  regras de inclusão e exclusão;
\item
  dimensões permitidas;
\item
  tabelas e colunas de origem;
\item
  interpretação e limitações.
\end{itemize}

Toda consulta financeira deve:

\begin{itemize}
\tightlist
\item
  declarar a granularidade final;
\item
  agregar pagamentos antes de combiná-los com itens;
\item
  contar pedidos distintos após joins 1:N;
\item
  impedir a repetição de valores de pedido por item ou pagamento;
\item
  declarar se utiliza todos os pedidos ou apenas entregues;
\item
  possuir consulta de controle;
\item
  separar medidas observadas de medidas alocadas.
\end{itemize}

\subsection{Rastreabilidade com os requisitos
aprovados}\label{docs__analytics__contrato_analitico.md__rastreabilidade-com-os-requisitos-aprovados}

{\def\LTcaptype{none} % do not increment counter
\begin{longtable}[]{@{}
  >{\raggedright\arraybackslash}p{(\linewidth - 4\tabcolsep) * \real{0.3333}}
  >{\raggedright\arraybackslash}p{(\linewidth - 4\tabcolsep) * \real{0.3333}}
  >{\raggedright\arraybackslash}p{(\linewidth - 4\tabcolsep) * \real{0.3333}}@{}}
\toprule\noalign{}
\begin{minipage}[b]{\linewidth}\raggedright
Demanda deste contrato
\end{minipage} & \begin{minipage}[b]{\linewidth}\raggedright
Requisito existente
\end{minipage} & \begin{minipage}[b]{\linewidth}\raggedright
Evidência futura
\end{minipage} \\
\midrule\noalign{}
\endhead
\bottomrule\noalign{}
\endlastfoot
Desempenho financeiro multidimensional & RF03 & Queries financeiras e
estruturas de consumo \\
Desempenho de produtos, categorias e vendedores & RF05 & Queries e
catálogo de vendedores \\
Formas, valores e condições de pagamento & RF06 & Análise de pagamentos
e parcelamento \\
Análises geográficas & RF09 & Recortes de origem e alcance \\
Consumo por BI & RF12 & Views, marts ou consultas de exportação
justificadas \\
Preservação de granularidade & Regra aprovada de joins analíticos &
Consultas de controle e validação \\
\end{longtable}
}

\subsection{Impacto no fluxo
M06}\label{docs__analytics__contrato_analitico.md__impacto-no-fluxo-m06}

\begin{itemize}
\tightlist
\item
  M06-01 deve explorar a CORE a partir das perguntas priorizadas neste
  contrato.
\item
  M06-02 deve combinar entidades para responder aos casos de uso
  aprovados.
\item
  M06-03 deve validar a cobertura dos requisitos e registrar evidências.
\item
  M06-04 deve validar e consolidar o catálogo preliminar de métricas.
\item
  M06-05 deve persistir somente regras reutilizáveis e semanticamente
  estáveis.
\item
  M06-06 deve preparar datasets com nomes, granularidade e métricas
  explícitos.
\end{itemize}

O contrato pode ser refinado durante essas entregas. Qualquer mudança
deve registrar a evidência, a decisão e o impacto sobre consultas,
métricas e estruturas já produzidas.

\section{Evidências das consultas básicas sobre a
CORE}\label{docs__analytics__evidencias_consultas_basicas_core.md__eviduxeancias-das-consultas-buxe1sicas-sobre-a-core}

\subsection{Finalidade}\label{docs__analytics__evidencias_consultas_basicas_core.md__finalidade}

Este documento registra os principais resultados de controle obtidos
durante a implementação das consultas básicas da etapa M06-01. Os
resultados foram produzidos sobre a CORE local validada e servem para
confirmar decisões de granularidade do contrato analítico.

As evidências não substituem os SQLs versionados em \nolinkurl{queries/}
e não transformam características observadas no dataset em regras de
negócio.

\subsection{Base
executada}\label{docs__analytics__evidencias_consultas_basicas_core.md__base-executada}

{\def\LTcaptype{none} % do not increment counter
\begin{longtable}[]{@{}lr@{}}
\toprule\noalign{}
Relação CORE & Registros \\
\midrule\noalign{}
\endhead
\bottomrule\noalign{}
\endlastfoot
\nolinkurl{core.pedido} & 99.441 \\
\nolinkurl{core.item_pedido} & 112.650 \\
\nolinkurl{core.pagamento} & 103.886 \\
\nolinkurl{core.vendedor} & 3.095 \\
\end{longtable}
}

As consultas foram executadas no PostgreSQL 18 contra o resultado
aprovado do ELT.

\subsection{Status dos
pedidos}\label{docs__analytics__evidencias_consultas_basicas_core.md__status-dos-pedidos}

{\def\LTcaptype{none} % do not increment counter
\begin{longtable}[]{@{}lrr@{}}
\toprule\noalign{}
Status & Pedidos & Participação \\
\midrule\noalign{}
\endhead
\bottomrule\noalign{}
\endlastfoot
delivered & 96.478 & 97,02\% \\
shipped & 1.107 & 1,11\% \\
canceled & 625 & 0,63\% \\
unavailable & 609 & 0,61\% \\
invoiced & 314 & 0,32\% \\
processing & 301 & 0,30\% \\
created & 5 & 0,01\% \\
approved & 2 & menor que 0,01\% \\
\end{longtable}
}

Decisão: os KPIs financeiros principais utilizarão pedidos entregues. Os
demais status permanecerão disponíveis para cobertura, reconciliação e
investigação.

\subsection{Cobertura
temporal}\label{docs__analytics__evidencias_consultas_basicas_core.md__cobertura-temporal}

\nolinkurl{core.pedido.data_compra} varia de setembro de 2016 a outubro
de 2018. Os meses de borda e transição que não possuem pedidos em todos
os dias são:

\begin{itemize}
\tightlist
\item
  setembro, outubro e dezembro de 2016;
\item
  janeiro de 2017;
\item
  setembro e outubro de 2018.
\end{itemize}

Fevereiro de 2017 a agosto de 2018 possuem ao menos um pedido em todos
os dias do respectivo mês. A ausência de pedidos em parte de um mês é um
sinal de cobertura parcial, não uma prova isolada de perda de dados.

Decisão: comparações mensais e taxas de crescimento devem destacar ou
excluir meses sem cobertura diária completa, conforme o objetivo da
análise.

\subsection{Cobertura das relações de
pedido}\label{docs__analytics__evidencias_consultas_basicas_core.md__cobertura-das-relauxe7uxf5es-de-pedido}

A maior parte dos pedidos possui itens e pagamentos. Foram observadas
exceções:

\begin{itemize}
\tightlist
\item
  um pedido entregue possui itens e não possui pagamento;
\item
  164 pedidos cancelados possuem pagamentos e não possuem itens;
\item
  603 pedidos indisponíveis possuem pagamentos e não possuem itens;
\item
  ocorrências menores também existem nos status \nolinkurl{created},
  \nolinkurl{invoiced} e \nolinkurl{shipped}.
\end{itemize}

Decisão: métricas que exigem itens e pagamentos devem utilizar
\nolinkurl{INNER JOIN} após pré-agregação e declarar que sua população
exclui pedidos sem uma das relações. Consultas de cobertura devem
preservar as exceções.

\subsection{Risco de multiplicação entre itens e
pagamentos}\label{docs__analytics__evidencias_consultas_basicas_core.md__risco-de-multiplicauxe7uxe3o-entre-itens-e-pagamentos}

O join direto entre \nolinkurl{core.item_pedido} e
\nolinkurl{core.pagamento} produziu 117.601 linhas. Sem pré-agregação,
itens são repetidos pela quantidade de pagamentos e pagamentos são
repetidos pela quantidade de itens do pedido.

A consulta de controle compara o join direto somente com pedidos
presentes nas duas relações. O resultado comprova que somas financeiras
realizadas após esse join não são válidas.

Decisão: itens e pagamentos devem ser agregados separadamente por
\nolinkurl{id_pedido} antes da combinação.

\subsection{Reconciliação
financeira}\label{docs__analytics__evidencias_consultas_basicas_core.md__reconciliauxe7uxe3o-financeira}

Entre os pedidos entregues que possuem itens e pagamentos:

\begin{itemize}
\tightlist
\item
  96.178 estão conciliados dentro da tolerância de R\$ 0,01;
\item
  260 possuem pagamento registrado maior que o valor bruto;
\item
  39 possuem valor bruto maior que o pagamento registrado;
\item
  a maior diferença absoluta observada é R\$ 182,81.
\end{itemize}

A tolerância absorve somente diferenças de centavos. Divergências
maiores não comprovam erro, estorno ou perda financeira, pois a fonte
não contém todo o ciclo contábil.

Decisão: a reconciliação será mantida como indicador de qualidade e
investigação, não como ajuste automático dos valores.

\subsection{Pedidos com múltiplos
vendedores}\label{docs__analytics__evidencias_consultas_basicas_core.md__pedidos-com-muxfaltiplos-vendedores}

Na base de pedidos entregues com itens:

{\def\LTcaptype{none} % do not increment counter
\begin{longtable}[]{@{}lrr@{}}
\toprule\noalign{}
Classificação & Pedidos & Participação \\
\midrule\noalign{}
\endhead
\bottomrule\noalign{}
\endlastfoot
Um vendedor & 95.203 & 98,6785\% \\
Múltiplos vendedores & 1.275 & 1,3215\% \\
\end{longtable}
}

Os pedidos multivendedor contêm 3.097 itens e R\$ 268.630,76 em valor
dos itens.

Decisão: pagamentos, avaliações, cancelamentos e atrasos de entrega não
serão atribuídos integralmente a cada vendedor. Quando usados, serão
descritos como indicadores associados ao pedido com participação do
vendedor.

\subsection{Reconciliação da atribuição aos
vendedores}\label{docs__analytics__evidencias_consultas_basicas_core.md__reconciliauxe7uxe3o-da-atribuiuxe7uxe3o-aos-vendedores}

Para pedidos entregues:

\begin{itemize}
\tightlist
\item
  110.197 itens da base foram atribuídos a vendedores;
\item
  a diferença entre itens da base e itens atribuídos foi zero;
\item
  R\$ 13.221.498,11 em valor dos itens foram atribuídos;
\item
  a diferença monetária da atribuição foi R\$ 0,00;
\item
  existem 1.341 participações adicionais em pedidos devido a pedidos com
  mais de um vendedor.
\end{itemize}

Decisão: valores e itens podem ser atribuídos diretamente aos
vendedores. A contagem de participações por vendedor não deve ser
confundida com pedidos distintos do marketplace.

\subsection{Consultas
relacionadas}\label{docs__analytics__evidencias_consultas_basicas_core.md__consultas-relacionadas}

\begin{itemize}
\tightlist
\item
  \nolinkurl{queries/basic/01_pedidos_por_status.sql};
\item
  \nolinkurl{queries/basic/02_cobertura_temporal.sql};
\item
  \nolinkurl{queries/basic/08_pedidos_multivendedor.sql};
\item
  \nolinkurl{queries/validation/01_reconciliacao_financeira.sql};
\item
  \nolinkurl{queries/validation/02_controle_join_itens_pagamentos.sql};
\item
  \nolinkurl{queries/validation/03_cobertura_relacoes_pedido.sql};
\item
  \nolinkurl{queries/validation/04_granularidade_vendedor.sql}.
\end{itemize}

\section{Evidências das consultas analíticas sobre a
CORE}\label{docs__analytics__evidencias_consultas_analiticas_core.md__eviduxeancias-das-consultas-analuxedticas-sobre-a-core}

\subsection{Escopo}\label{docs__analytics__evidencias_consultas_analiticas_core.md__escopo}

Este documento registra evidências da etapa M06-02, orientada
prioritariamente ao desempenho financeiro e à contribuição dos
vendedores. As consultas-fonte estão em \nolinkurl{queries/analysis/} e
utilizam pedidos entregues como base principal.

Os resultados descrevem o período histórico do dataset. Não representam
a operação atual, receita contábil, lucro ou participação no mercado
brasileiro.

\subsection{Resumo financeiro
histórico}\label{docs__analytics__evidencias_consultas_analiticas_core.md__resumo-financeiro-histuxf3rico}

A população que possui simultaneamente pedido entregue, itens e
pagamentos contém 96.477 pedidos e 110.194 itens.

{\def\LTcaptype{none} % do not increment counter
\begin{longtable}[]{@{}lr@{}}
\toprule\noalign{}
Indicador & Resultado \\
\midrule\noalign{}
\endhead
\bottomrule\noalign{}
\endlastfoot
Valor dos itens & R\$ 13.221.363,14 \\
Valor de frete & R\$ 2.198.267,15 \\
Valor bruto & R\$ 15.419.630,29 \\
Valor pago registrado & R\$ 15.422.461,77 \\
Diferença agregada de reconciliação & R\$ 2.831,48 \\
Valor bruto médio por pedido & R\$ 159,83 \\
Participação do frete & 14,26\% \\
\end{longtable}
}

A diferença agregada não significa que todos os pedidos estejam
conciliados. Diferenças positivas e negativas podem se compensar; a
validação por pedido permanece necessária.

\subsection{Evolução
mensal}\label{docs__analytics__evidencias_consultas_analiticas_core.md__evoluuxe7uxe3o-mensal}

A consulta mensal:

\begin{itemize}
\tightlist
\item
  utiliza \nolinkurl{core.pedido.data_compra};
\item
  agrega itens e pagamentos separadamente por pedido;
\item
  sinaliza meses sem pedidos em todos os dias do calendário;
\item
  calcula crescimento somente quando o mês atual e o anterior possuem
  cobertura diária completa.
\end{itemize}

Essa regra evita comparar automaticamente os meses de borda ou transição
já identificados na etapa M06-01.

\subsection{Concentração dos
vendedores}\label{docs__analytics__evidencias_consultas_analiticas_core.md__concentrauxe7uxe3o-dos-vendedores}

Foram observados 2.970 vendedores ativos em pedidos entregues.

{\def\LTcaptype{none} % do not increment counter
\begin{longtable}[]{@{}lr@{}}
\toprule\noalign{}
Indicador & Resultado \\
\midrule\noalign{}
\endhead
\bottomrule\noalign{}
\endlastfoot
Participação acumulada dos 5 maiores & 7,71\% \\
Participação acumulada dos 10 maiores & 13,27\% \\
HHI da base observada & 0,003631 \\
\end{longtable}
}

Os resultados indicam dispersão do valor dos itens entre muitos
vendedores no período completo. O HHI é utilizado somente como medida
descritiva interna. Não permite concluir sobre concentração
concorrencial ou participação de mercado.

A curva ABC usa inicialmente os limites acumulados de 80\% e 95\%. Esses
limites são parâmetros analíticos e poderão ser revistos na definição
final das métricas.

\subsection{Pedidos com múltiplos
vendedores}\label{docs__analytics__evidencias_consultas_analiticas_core.md__pedidos-com-muxfaltiplos-vendedores}

Entre os pedidos entregues com itens e pagamentos, 1.275 possuem mais de
um vendedor. Por isso:

\begin{itemize}
\tightlist
\item
  valor dos itens e frete podem ser atribuídos diretamente ao vendedor;
\item
  pagamentos não são repetidos nem integralmente atribuídos aos
  vendedores;
\item
  atrasos e avaliações são descritos como associados ao pedido com
  participação do vendedor.
\end{itemize}

\subsection{Formas e condições de
pagamento}\label{docs__analytics__evidencias_consultas_analiticas_core.md__formas-e-condiuxe7uxf5es-de-pagamento}

{\def\LTcaptype{none} % do not increment counter
\begin{longtable}[]{@{}
  >{\raggedright\arraybackslash}p{(\linewidth - 8\tabcolsep) * \real{0.1579}}
  >{\raggedleft\arraybackslash}p{(\linewidth - 8\tabcolsep) * \real{0.2105}}
  >{\raggedleft\arraybackslash}p{(\linewidth - 8\tabcolsep) * \real{0.2105}}
  >{\raggedleft\arraybackslash}p{(\linewidth - 8\tabcolsep) * \real{0.2105}}
  >{\raggedleft\arraybackslash}p{(\linewidth - 8\tabcolsep) * \real{0.2105}}@{}}
\toprule\noalign{}
\begin{minipage}[b]{\linewidth}\raggedright
Tipo
\end{minipage} & \begin{minipage}[b]{\linewidth}\raggedleft
Registros
\end{minipage} & \begin{minipage}[b]{\linewidth}\raggedleft
Pedidos distintos
\end{minipage} & \begin{minipage}[b]{\linewidth}\raggedleft
Valor pago registrado
\end{minipage} & \begin{minipage}[b]{\linewidth}\raggedleft
Parcelas médias
\end{minipage} \\
\midrule\noalign{}
\endhead
\bottomrule\noalign{}
\endlastfoot
Cartão de crédito & 74.586 & 74.304 & R\$ 12.101.094,88 & 3,50 \\
Boleto & 19.191 & 19.191 & R\$ 2.769.932,58 & 1,00 \\
Voucher & 5.493 & 3.679 & R\$ 343.013,19 & 1,00 \\
Cartão de débito & 1.486 & 1.485 & R\$ 208.421,12 & 1,00 \\
\end{longtable}
}

Um pedido pode possuir vários registros e tipos de pagamento. As
contagens por tipo não formam grupos mutuamente exclusivos de pedidos. A
consulta de perfil consolida as sequências por pedido antes de
classificá-lo por faixa de valor, parcelamento e composição dos meios de
pagamento.

\subsection{Risco operacional associado aos
vendedores}\label{docs__analytics__evidencias_consultas_analiticas_core.md__risco-operacional-associado-aos-vendedores}

Na granularidade vendedor-pedido foram observadas 97.819 participações.

{\def\LTcaptype{none} % do not increment counter
\begin{longtable}[]{@{}
  >{\raggedright\arraybackslash}p{(\linewidth - 2\tabcolsep) * \real{0.4286}}
  >{\raggedleft\arraybackslash}p{(\linewidth - 2\tabcolsep) * \real{0.5714}}@{}}
\toprule\noalign{}
\begin{minipage}[b]{\linewidth}\raggedright
Indicador associado
\end{minipage} & \begin{minipage}[b]{\linewidth}\raggedleft
Resultado
\end{minipage} \\
\midrule\noalign{}
\endhead
\bottomrule\noalign{}
\endlastfoot
Participação vendedor-pedido com algum item enviado após o limite &
8,97\% \\
Pedido entregue após a estimativa & 8,02\% \\
Nota média & 4,14 \\
Avaliação negativa, nota menor ou igual a 2 & 13,25\% \\
\end{longtable}
}

As taxas usam somente observações com informação disponível no
denominador. O resultado de 8,97\% possui granularidade vendedor-pedido
e indica a presença de ao menos um item fora do limite. A taxa calculada
na granularidade de item é 9,32\%, conforme estabilizado na M06-05.
Envio após o limite é diretamente relacionado aos itens do vendedor.
Atraso da entrega e avaliação pertencem ao pedido e não comprovam
responsabilidade nem causalidade do vendedor.

\subsection{Recortes de categoria e
geografia}\label{docs__analytics__evidencias_consultas_analiticas_core.md__recortes-de-categoria-e-geografia}

As consultas de vendedor-categoria e alcance geográfico permitem:

\begin{itemize}
\tightlist
\item
  medir a participação de uma categoria no valor do vendedor;
\item
  medir a participação do vendedor dentro da categoria;
\item
  contar estados e cidades de clientes atendidos;
\item
  comparar o peso do frete no valor bruto por vendedor.
\end{itemize}

O estado e a cidade do cliente representam seu cadastro, não o endereço
exato da entrega. Alcance geográfico não equivale a cobertura
operacional formal.

\subsection{Decisões para as próximas
etapas}\label{docs__analytics__evidencias_consultas_analiticas_core.md__decisuxf5es-para-as-pruxf3ximas-etapas}

\begin{itemize}
\tightlist
\item
  manter as bases de item, pagamento e avaliação pré-agregadas;
\item
  preservar meses parciais, mas impedir seu uso automático em
  crescimento;
\item
  tratar concentração como característica interna da base;
\item
  separar medidas diretamente atribuíveis de indicadores apenas
  associados;
\item
  validar fórmulas e cobertura contra os requisitos na M06-03;
\item
  consolidar como métricas somente os indicadores aprovados na M06-04.
\end{itemize}

\section{Matriz de validação dos requisitos
funcionais}\label{docs__analytics__matriz_validacao_requisitos_funcionais.md__matriz-de-validauxe7uxe3o-dos-requisitos-funcionais}

\subsection{Finalidade}\label{docs__analytics__matriz_validacao_requisitos_funcionais.md__finalidade}

Esta matriz registra como os requisitos funcionais RF01 a RF12 são
sustentados pelo modelo CORE e pelas consultas versionadas nas etapas
M06-01 e M06-02. A validação comprova capacidade analítica e cobertura
observada; não declara que toda análise possível já foi implementada nem
antecipa as estruturas de BI das etapas posteriores.

Classificações utilizadas:

\begin{itemize}
\tightlist
\item
  \textbf{atendido}: há estrutura, consulta e resultado reproduzível
  suficientes para demonstrar a capacidade;
\item
  \textbf{atendido com limitação}: a capacidade existe, mas a cobertura
  ou a interpretação é restringida pelos dados;
\item
  \textbf{preparado para etapa futura}: a base está disponível, mas a
  entrega final pertence explicitamente a uma issue posterior.
\end{itemize}

\subsection{Matriz requisito → consulta →
evidência}\label{docs__analytics__matriz_validacao_requisitos_funcionais.md__matriz-requisito-consulta-eviduxeancia}

{\def\LTcaptype{none} % do not increment counter
\begin{longtable}[]{@{}
  >{\raggedright\arraybackslash}p{(\linewidth - 8\tabcolsep) * \real{0.2000}}
  >{\raggedright\arraybackslash}p{(\linewidth - 8\tabcolsep) * \real{0.2000}}
  >{\raggedright\arraybackslash}p{(\linewidth - 8\tabcolsep) * \real{0.2000}}
  >{\raggedright\arraybackslash}p{(\linewidth - 8\tabcolsep) * \real{0.2000}}
  >{\raggedright\arraybackslash}p{(\linewidth - 8\tabcolsep) * \real{0.2000}}@{}}
\toprule\noalign{}
\begin{minipage}[b]{\linewidth}\raggedright
Requisito
\end{minipage} & \begin{minipage}[b]{\linewidth}\raggedright
Capacidade validada
\end{minipage} & \begin{minipage}[b]{\linewidth}\raggedright
Consulta ou artefato
\end{minipage} & \begin{minipage}[b]{\linewidth}\raggedright
Evidência observada
\end{minipage} & \begin{minipage}[b]{\linewidth}\raggedright
Resultado
\end{minipage} \\
\midrule\noalign{}
\endhead
\bottomrule\noalign{}
\endlastfoot
RF01 & Integração de clientes, pedidos, itens, produtos, vendedores,
pagamentos, avaliações e geolocalização &
\nolinkurl{queries/validation/05_cobertura_requisitos_funcionais.sql},
bloco RF01 & Pedidos se relacionam integralmente a clientes e itens a
produtos e vendedores. Pagamentos, avaliações e geolocalização possuem
cobertura parcial da fonte, preservada sem fabricação de registros. &
Atendido com limitação \\
RF02 & Recuperação do histórico de pedidos e relações &
\nolinkurl{queries/basic/01_pedidos_por_status.sql},
\nolinkurl{queries/validation/03_cobertura_relacoes_pedido.sql} e bloco
RF02/RF04 da consulta 05 & 99.441 pedidos entre 2016-09-04 e 2018-10-17;
exceções de itens e pagamentos permanecem identificáveis por status. &
Atendido \\
RF03 & Análise multidimensional por período, produto, categoria,
cliente, vendedor e localização &
\nolinkurl{queries/basic/09_valor_por_categoria_regiao.sql},
\nolinkurl{queries/analysis/01_evolucao_financeira_mensal.sql},
\nolinkurl{03_desempenho_vendedor_categoria.sql} e
\nolinkurl{04_alcance_geografico_vendedores.sql} & Resultados
reproduzíveis nas granularidades mês, categoria-estado,
vendedor-categoria e vendedor, com dimensões ligadas ao pedido. &
Atendido \\
RF04 & Comportamento de compra a partir do histórico disponível & Bloco
RF02/RF04 de
\nolinkurl{queries/validation/05_cobertura_requisitos_funcionais.sql} &
O identificador persistente permite contar recorrência, intervalo de
compras e máximo de pedidos por cliente observado. Não há navegação,
abandono ou conversão. & Atendido com limitação \\
RF05 & Desempenho comercial de produtos, categorias e vendedores &
\nolinkurl{queries/basic/09_valor_por_categoria_regiao.sql},
\nolinkurl{queries/analysis/02_concentracao_vendedores.sql} e
\nolinkurl{03_desempenho_vendedor_categoria.sql} & Valor, itens,
pedidos, participação, ranking, concentração e recortes por categoria
são calculáveis sem atribuir pagamentos ao vendedor. & Atendido \\
RF06 & Formas, valores e condições de pagamento &
\nolinkurl{queries/basic/06_composicao_pagamentos.sql} e
\nolinkurl{queries/analysis/06_perfil_pagamentos_pedido.sql} & Tipos,
valores, parcelas e composição de meios de pagamento são preservados;
pedidos com múltiplas sequências são agregados antes da análise. &
Atendido \\
RF07 & Relação entre avaliações e demais informações do pedido &
\nolinkurl{queries/analysis/05_risco_operacional_vendedores.sql} e bloco
RF07/RF08 da consulta 05 & Avaliações podem ser pré-agregadas por pedido
e associadas a itens e vendedores. A associação não comprova
responsabilidade ou causalidade. & Atendido com limitação \\
RF08 & Análise de datas previstas e efetivas de processamento e entrega
& \nolinkurl{queries/analysis/05_risco_operacional_vendedores.sql} e
bloco RF07/RF08 da consulta 05 & Aprovação, envio, entrega, estimativa e
envio após limite permitem medir cobertura e atrasos quando as datas
estão disponíveis. & Atendido com limitação \\
RF09 & Análises geográficas de clientes e vendedores &
\nolinkurl{queries/analysis/04_alcance_geografico_vendedores.sql},
\nolinkurl{queries/basic/09_valor_por_categoria_regiao.sql} e bloco RF09
da consulta 05 & Estados e cidades permitem recortes comerciais;
prefixos podem ser relacionados à geolocalização quando cobertos.
Coordenadas representam aproximações do prefixo, não rotas. & Atendido
com limitação \\
RF10 & Combinação de entidades com integridade dos relacionamentos &
\nolinkurl{queries/validation/01_reconciliacao_financeira.sql},
\nolinkurl{02_controle_join_itens_pagamentos.sql} e
\nolinkurl{04_granularidade_vendedor.sql} & A pré-agregação por pedido
evita multiplicar itens e pagamentos; a atribuição de itens e valor aos
vendedores reconciliou sem diferença. & Atendido \\
RF11 & Agregações e filtros temporais, comerciais, geográficos e de
pedido & Coleções \nolinkurl{queries/basic/} e
\nolinkurl{queries/analysis/} & Consultas filtram status, controlam
meses parciais e agregam por mês, categoria, região, vendedor, pagamento
e pedido. & Atendido \\
RF12 & Base estruturada para ferramentas analíticas e BI & Schemas
\nolinkurl{core} e \nolinkurl{analytics}, contrato analítico e consultas
versionadas & A CORE validada e as regras reproduzíveis formam a base de
consumo. Views, marts e datasets finais serão materializados nas issues
M06-05 e M06-06. & Preparado para etapa futura \\
\end{longtable}
}

\subsection{Evidências complementares da consulta de
cobertura}\label{docs__analytics__matriz_validacao_requisitos_funcionais.md__eviduxeancias-complementares-da-consulta-de-cobertura}

\subsubsection{Integração e
cobertura}\label{docs__analytics__matriz_validacao_requisitos_funcionais.md__integrauxe7uxe3o-e-cobertura}

As chaves estrangeiras asseguram as relações obrigatórias. Relações
opcionais ou incompletas refletem a fonte: a ausência de pagamento,
avaliação ou geolocalização não deve ser preenchida artificialmente. A
consulta registra a quantidade relacionada e não relacionada para tornar
essa cobertura auditável.

{\def\LTcaptype{none} % do not increment counter
\begin{longtable}[]{@{}lrrr@{}}
\toprule\noalign{}
Relação validada & Origem & Relacionados & Sem relação \\
\midrule\noalign{}
\endhead
\bottomrule\noalign{}
\endlastfoot
Item → produto e vendedor & 112.650 & 112.650 & 0 \\
Pedido → cliente & 99.441 & 99.441 & 0 \\
Pedido → pagamento & 99.441 & 99.440 & 1 \\
Pedido → avaliação & 99.441 & 98.673 & 768 \\
Prefixo de CEP → geolocalização & 19.177 & 19.015 & 162 \\
\end{longtable}
}

\subsubsection{Histórico e comportamento de
clientes}\label{docs__analytics__matriz_validacao_requisitos_funcionais.md__histuxf3rico-e-comportamento-de-clientes}

O histórico usa \nolinkurl{core.cliente.id_cliente_unico} como
identidade persistente e \nolinkurl{core.pedido.data_compra} como
referência temporal. A análise fica limitada a clientes com transações
observadas; não existem dados de sessão, visita, abandono ou exposição
comercial.

Foram observados 96.096 clientes persistentes; 2.997 (3,12\%) possuem
mais de um pedido, e o máximo observado é 17 pedidos por cliente. O
histórico varia de 2016-09-04 a 2018-10-17.

\subsubsection{Avaliação, entrega e
geografia}\label{docs__analytics__matriz_validacao_requisitos_funcionais.md__avaliauxe7uxe3o-entrega-e-geografia}

Avaliações e datas são associadas ao pedido somente quando disponíveis.
Cidade, estado e prefixo de CEP representam os atributos cadastrais
presentes na base. As múltiplas ocorrências de geolocalização por
prefixo não devem ser ligadas diretamente a fatos sem agregação prévia.

Entre os 99.441 pedidos, 98.673 possuem avaliação, 99.281 possuem data
de aprovação, 97.658 possuem envio ao transportador, 96.476 possuem
entrega e todos possuem estimativa. Foram observados 7.827 pedidos
entregues após a data estimada, sem atribuição causal.

A geolocalização cobre 99.163 de 99.441 registros de cliente e 3.088 de
3.095 vendedores. Os cadastros abrangem 27 estados e 4.119 cidades de
clientes, além de 23 estados e 611 cidades de vendedores.

\subsection{Decisões e impacto nas próximas
etapas}\label{docs__analytics__matriz_validacao_requisitos_funcionais.md__decisuxf5es-e-impacto-nas-pruxf3ximas-etapas}

\begin{itemize}
\tightlist
\item
  RF01 a RF11 possuem cobertura suficiente para prosseguir à definição
  formal das métricas, respeitadas as limitações registradas.
\item
  RF04 não autoriza métricas de conversão, abandono ou navegação.
\item
  RF07 e RF08 sustentam indicadores associados ao vendedor, mas não
  atribuição causal de avaliações ou atraso final.
\item
  RF09 não sustenta distância percorrida nem cobertura logística formal.
\item
  RF12 somente será classificado como entrega final após a criação das
  estruturas reutilizáveis e dos datasets de consumo das issues M06-05 e
  M06-06.
\item
  Nenhuma evidência exige alterar os modelos conceitual, lógico, físico
  ou a transformação CORE já aprovados.
\end{itemize}

\subsection{Reprodutibilidade}\label{docs__analytics__matriz_validacao_requisitos_funcionais.md__reprodutibilidade}

As evidências devem ser reproduzidas no PostgreSQL 18 sobre a CORE
aprovada, executando as consultas referenciadas nesta matriz. As
contagens documentadas correspondem à carga local validada e podem
variar somente se a fonte ou o ELT forem formalmente alterados.

\section{Catálogo de métricas
analíticas}\label{docs__analytics__catalogo_metricas_analiticas.md__catuxe1logo-de-muxe9tricas-analuxedticas}

\subsection{Finalidade e
escopo}\label{docs__analytics__catalogo_metricas_analiticas.md__finalidade-e-escopo}

Este catálogo consolida as métricas aprovadas para responder
prioritariamente ao desempenho financeiro e, em seguida, ao desempenho
dos vendedores. As definições utilizam exclusivamente a CORE validada e
correspondem ao período histórico do dataset.

Os valores observados não representam receita contábil, lucro, margem,
liquidação financeira nem a operação atual da Olist. Todas as métricas
são descritivas e não demonstram causalidade.

\subsection{Políticas
comuns}\label{docs__analytics__catalogo_metricas_analiticas.md__poluxedticas-comuns}

\begin{itemize}
\tightlist
\item
  \textbf{Data financeira e comercial:}
  \nolinkurl{core.pedido.data_compra}.
\item
  \textbf{Base principal:} pedidos com
  \nolinkurl{status_pedido = 'delivered'}.
\item
  \textbf{Base financeira completa:} pedidos entregues com itens e
  pagamentos, quando a métrica compara essas duas representações.
\item
  \textbf{Cancelamento:} todos os pedidos compõem o denominador; o
  numerador contém somente \nolinkurl{status_pedido = 'canceled'}.
\item
  \textbf{Granularidade:} itens e pagamentos são agregados separadamente
  por pedido antes de qualquer combinação.
\item
  \textbf{Vendedores:} valor dos itens, frete e quantidade de itens são
  diretamente atribuíveis. Entrega final e avaliação são apenas
  associadas ao pedido com participação do vendedor.
\item
  \textbf{Percentuais:} são expressos na escala de 0 a 100. HHI
  permanece na escala de 0 a 1.
\item
  \textbf{Valores monetários:} BRL com precisão exata; arredondamento é
  aplicado somente à apresentação de médias, taxas e percentuais.
\end{itemize}

\subsection{Métricas
financeiras}\label{docs__analytics__catalogo_metricas_analiticas.md__muxe9tricas-financeiras}

{\def\LTcaptype{none} % do not increment counter
\begin{longtable}[]{@{}
  >{\raggedright\arraybackslash}p{(\linewidth - 12\tabcolsep) * \real{0.1429}}
  >{\raggedright\arraybackslash}p{(\linewidth - 12\tabcolsep) * \real{0.1429}}
  >{\raggedright\arraybackslash}p{(\linewidth - 12\tabcolsep) * \real{0.1429}}
  >{\raggedright\arraybackslash}p{(\linewidth - 12\tabcolsep) * \real{0.1429}}
  >{\raggedright\arraybackslash}p{(\linewidth - 12\tabcolsep) * \real{0.1429}}
  >{\raggedright\arraybackslash}p{(\linewidth - 12\tabcolsep) * \real{0.1429}}
  >{\raggedright\arraybackslash}p{(\linewidth - 12\tabcolsep) * \real{0.1429}}@{}}
\toprule\noalign{}
\begin{minipage}[b]{\linewidth}\raggedright
ID
\end{minipage} & \begin{minipage}[b]{\linewidth}\raggedright
Nome e finalidade
\end{minipage} & \begin{minipage}[b]{\linewidth}\raggedright
Fórmula e unidade
\end{minipage} & \begin{minipage}[b]{\linewidth}\raggedright
Granularidade e tempo
\end{minipage} & \begin{minipage}[b]{\linewidth}\raggedright
Filtros e exclusões
\end{minipage} & \begin{minipage}[b]{\linewidth}\raggedright
Fontes e SQL
\end{minipage} & \begin{minipage}[b]{\linewidth}\raggedright
Interpretação e limitações
\end{minipage} \\
\midrule\noalign{}
\endhead
\bottomrule\noalign{}
\endlastfoot
FIN01 & \textbf{Valor dos itens entregues} --- medir o valor comercial
dos produtos & \nolinkurl{SUM(preco_item)}; BRL & Origem item; saída
histórica, mensal ou segmento; \nolinkurl{data_compra} & Pedidos
entregues; na comparação financeira, exige item e pagamento &
\nolinkurl{core.item_pedido.preco_item}; consultas 01 e 07 de
\nolinkurl{queries/analysis/} & Valor observado dos itens; não é receita
líquida, margem ou lucro. \\
FIN02 & \textbf{Valor de frete entregue} --- medir o frete cobrado nos
itens & \nolinkurl{SUM(valor_frete)}; BRL & Origem item; saída
histórica, mensal ou segmento; \nolinkurl{data_compra} & Mesma base da
FIN01 & \nolinkurl{core.item_pedido.valor_frete}; consultas 01 e 07 &
Frete cobrado, não custo logístico real. \\
FIN03 & \textbf{Valor bruto entregue} --- medir itens mais frete &
\nolinkurl{SUM(preco_item + valor_frete)}; BRL & Item, pré-agregado por
pedido; saída histórica, mensal ou segmento & Mesma base da FIN01 &
\nolinkurl{core.item_pedido}; consultas 01 e 07 & Valor bruto observado;
não comprova reconhecimento contábil. \\
FIN04 & \textbf{Valor pago registrado} --- medir pagamentos associados
aos pedidos & \nolinkurl{SUM(valor_pagamento)}; BRL & Sequência de
pagamento, pré-agregada por pedido; saída histórica ou mensal & Pedidos
entregues com itens e pagamentos &
\nolinkurl{core.pagamento.valor_pagamento}; consultas 01 e 07 & Registro
da fonte; não comprova liquidação, estorno ou recebimento. \\
FIN05 & \textbf{Diferença de reconciliação} --- comparar pagamento e
valor bruto & \nolinkurl{valor_pago_registrado - valor_bruto}; BRL por
pedido, com soma e classificação & Origem e controle por pedido; saída
mensal ou histórica; \nolinkurl{data_compra} & Somente pedidos presentes
em itens e pagamentos; tolerância de R\$ 0,01 apenas na classificação &
\nolinkurl{core.item_pedido}, \nolinkurl{core.pagamento}; análise 01 e
validação 01 & Divergência para investigação; somas positivas e
negativas podem se compensar. \\
FIN06 & \textbf{Valor bruto médio por pedido entregue} --- medir valor
bruto médio da base &
\nolinkurl{SUM(valor_bruto) / COUNT(DISTINCT id_pedido)}; BRL/pedido &
Pedido; saída histórica, mensal ou segmento; \nolinkurl{data_compra} &
Pedidos entregues com itens e pagamentos & Análises 01 e 07 & Não
utilizar o nome ambíguo ``ticket médio''; a base monetária é
explícita. \\
FIN07 & \textbf{Participação do frete} --- medir o peso do frete no
valor bruto & \nolinkurl{100 × SUM(valor_frete) / SUM(valor_bruto)};
percentual & Saída histórica, mensal ou segmento;
\nolinkurl{data_compra} & Pedidos entregues; denominador diferente de
zero & Análises 01 e 07 & Mede composição do valor cobrado, não
eficiência ou custo logístico. \\
FIN08 & \textbf{Crescimento mensal do valor dos itens} --- comparar
evolução mensal &
\nolinkurl{100 × (valor_itens_mês / valor_itens_mês_anterior - 1)};
percentual & Mês de compra & Mês atual e anterior devem possuir
cobertura diária completa; pedidos entregues com itens e pagamentos &
\nolinkurl{queries/analysis/01_evolucao_financeira_mensal.sql} & Meses
parciais retornam nulo; variação não explica causa. \\
FIN09 & \textbf{Taxa de cancelamento por quantidade} --- medir
participação de pedidos cancelados &
\nolinkurl{100 × pedidos_cancelados / todos_pedidos}; percentual &
Pedido, com saída mensal ou histórica; \nolinkurl{data_compra} & Todos
os status no denominador; \nolinkurl{canceled} no numerador &
\nolinkurl{core.pedido.status_pedido}; análise 08 & Mede registros
cancelados na base, não probabilidade futura. Meses parciais devem ser
sinalizados. \\
FIN10 & \textbf{Valores associados a pedidos cancelados} --- dimensionar
registros financeiros presentes & Somas separadas de itens, frete, bruto
e pagamento; BRL & Itens e pagamentos pré-agregados por pedido; saída
mensal ou histórica & Somente pedidos cancelados; ausências preservadas
como cobertura &
\nolinkurl{queries/analysis/08_metricas_cancelamento.sql} & Não
representa perda, estorno ou receita cancelada; as representações não
devem ser fundidas. \\
\end{longtable}
}

\subsection{Métricas de
vendedores}\label{docs__analytics__catalogo_metricas_analiticas.md__muxe9tricas-de-vendedores}

{\def\LTcaptype{none} % do not increment counter
\begin{longtable}[]{@{}
  >{\raggedright\arraybackslash}p{(\linewidth - 12\tabcolsep) * \real{0.1429}}
  >{\raggedright\arraybackslash}p{(\linewidth - 12\tabcolsep) * \real{0.1429}}
  >{\raggedright\arraybackslash}p{(\linewidth - 12\tabcolsep) * \real{0.1429}}
  >{\raggedright\arraybackslash}p{(\linewidth - 12\tabcolsep) * \real{0.1429}}
  >{\raggedright\arraybackslash}p{(\linewidth - 12\tabcolsep) * \real{0.1429}}
  >{\raggedright\arraybackslash}p{(\linewidth - 12\tabcolsep) * \real{0.1429}}
  >{\raggedright\arraybackslash}p{(\linewidth - 12\tabcolsep) * \real{0.1429}}@{}}
\toprule\noalign{}
\begin{minipage}[b]{\linewidth}\raggedright
ID
\end{minipage} & \begin{minipage}[b]{\linewidth}\raggedright
Nome e finalidade
\end{minipage} & \begin{minipage}[b]{\linewidth}\raggedright
Fórmula e unidade
\end{minipage} & \begin{minipage}[b]{\linewidth}\raggedright
Granularidade e tempo
\end{minipage} & \begin{minipage}[b]{\linewidth}\raggedright
Filtros e exclusões
\end{minipage} & \begin{minipage}[b]{\linewidth}\raggedright
Fontes e SQL
\end{minipage} & \begin{minipage}[b]{\linewidth}\raggedright
Interpretação e limitações
\end{minipage} \\
\midrule\noalign{}
\endhead
\bottomrule\noalign{}
\endlastfoot
VEN01 & \textbf{Vendedores ativos} --- medir a base com atividade
observada & \nolinkurl{COUNT(DISTINCT id_vendedor)}; vendedores & Item,
saída histórica ou por período; \nolinkurl{data_compra} & Ao menos um
item em pedido entregue no período & \nolinkurl{core.item_pedido},
\nolinkurl{core.pedido}; análises 02 e 07 & Atividade no dataset, não
situação contratual ou operacional atual. \\
VEN02 & \textbf{Valor dos itens por vendedor} --- medir contribuição
comercial direta & \nolinkurl{SUM(preco_item)}; BRL/vendedor & Vendedor;
histórica, período ou segmento; \nolinkurl{data_compra} & Itens de
pedidos entregues & Análise 02 e consulta básica 07 & Diretamente
atribuível ao vendedor; não é receita, lucro ou pagamento recebido. \\
VEN03 & \textbf{Participação financeira do vendedor} --- medir
relevância relativa &
\nolinkurl{100 × valor_itens_vendedor / valor_itens_total}; percentual &
Vendedor no mesmo período e população & Mesma base da VEN02; denominador
diferente de zero &
\nolinkurl{queries/analysis/02_concentracao_vendedores.sql} &
Participação na base observada, não participação de mercado. \\
VEN04 & \textbf{Participação acumulada dos maiores vendedores} --- medir
concentração & Soma das participações ordenadas por valor; percentual
Top 5, 10, 20 e 10\% da base & Ranking de vendedor por período & Mesma
base da VEN02; desempate por \nolinkurl{id_vendedor} & Análises 02 e 07
& Curva descritiva; classes ABC usam inicialmente 80\% e 95\%. \\
VEN05 & \textbf{Pedidos com participação do vendedor} --- medir alcance
em pedidos & \nolinkurl{COUNT(DISTINCT id_pedido)}; pedidos/vendedor &
Vendedor; histórica ou por período; \nolinkurl{data_compra} & Pedidos
entregues com item do vendedor & Consulta básica 07 & Em pedido
multivendedor, cada vendedor recebe uma participação; a soma excede
pedidos do marketplace. \\
VEN06 & \textbf{Itens por vendedor} --- medir volume comercializado &
\nolinkurl{COUNT(*)}; itens/vendedor & Vendedor; histórica ou por
período; \nolinkurl{data_compra} & Itens de pedidos entregues & Consulta
básica 07 & Mede linhas de item, não quantidade física além da
representação da fonte. \\
VEN07 & \textbf{Valor médio do vendedor por pedido com participação} ---
medir contribuição média &
\nolinkurl{valor_itens_vendedor / pedidos_distintos_com_participação};
BRL/pedido participante & Vendedor; histórica ou por período;
\nolinkurl{data_compra} & Mesma base da VEN02 & Consulta básica 07 &
Considera somente a parcela de itens do vendedor; não é o valor total
nem o pagamento do pedido. \\
VEN08 & \textbf{Taxa de itens enviados após o limite} --- medir
cumprimento do prazo diretamente relacionado ao item &
\nolinkurl{100 × itens após data_limite_envio / itens com observação};
percentual & Vendedor-item, consolidada por vendedor;
\nolinkurl{data_compra} & Pedidos entregues; somente observações não
nulas & \nolinkurl{core.item_pedido.data_limite_envio},
\nolinkurl{core.pedido.data_envio_transportador};
\nolinkurl{analytics.vw_desempenho_vendedor} & O envio é comparado ao
limite de cada item; não confundir com a proporção de pares
vendedor-pedido que possuem algum item fora do limite. \\
VEN09 & \textbf{Taxa de pedidos atrasados associados ao vendedor} ---
contextualizar risco de entrega &
\nolinkurl{100 × pedidos entregues após estimativa / pedidos com datas};
percentual & Vendedor-pedido; histórica ou por período;
\nolinkurl{data_compra} & Pedidos entregues com datas de entrega e
estimativa & \nolinkurl{core.pedido};
\nolinkurl{queries/analysis/05_risco_operacional_vendedores.sql} &
Atraso pertence ao pedido; usar ``associado'', sem atribuir
responsabilidade ao vendedor. \\
VEN10 & \textbf{Nota média associada ao vendedor} --- contextualizar
experiência & Média da nota média pré-agregada por pedido; pontos de 1 a
5 & Vendedor-pedido; histórica ou por período; \nolinkurl{data_compra} &
Pedidos entregues com avaliação &
\nolinkurl{core.avaliacao.nota_avaliacao}; análise 05 & Avaliação
pertence ao pedido e pode envolver vários vendedores; não mede
causalidade. \\
\end{longtable}
}

\subsection{Resultados históricos de
referência}\label{docs__analytics__catalogo_metricas_analiticas.md__resultados-histuxf3ricos-de-referuxeancia}

Na base de 96.477 pedidos entregues com itens e pagamentos:

{\def\LTcaptype{none} % do not increment counter
\begin{longtable}[]{@{}lr@{}}
\toprule\noalign{}
Métrica & Resultado \\
\midrule\noalign{}
\endhead
\bottomrule\noalign{}
\endlastfoot
FIN01 --- valor dos itens & R\$ 13.221.363,14 \\
FIN02 --- valor de frete & R\$ 2.198.267,15 \\
FIN03 --- valor bruto & R\$ 15.419.630,29 \\
FIN04 --- valor pago registrado & R\$ 15.422.461,77 \\
FIN05 --- diferença agregada & R\$ 2.831,48 \\
FIN06 --- valor bruto médio por pedido & R\$ 159,83 \\
FIN07 --- participação do frete & 14,26\% \\
\end{longtable}
}

No histórico completo, 625 de 99.441 pedidos estão cancelados (FIN09:
0,6285\%). Para esses pedidos, foram observados R\$ 95.235,27 em itens,
R\$ 10.650,45 em frete, R\$ 105.885,72 em valor bruto e R\$ 143.255,60
em pagamentos registrados. Há 164 cancelados sem itens e nenhum sem
pagamento, o que impede tratar as representações como equivalentes.

Para vendedores, foram observados 2.970 ativos. Os 5, 10 e 20 maiores
representam, respectivamente, 7,71\%, 13,27\% e 21,28\% do valor dos
itens; os 10\% maiores representam 67,11\%. Na granularidade de item, a
taxa de envio após o limite é 9,32\%. Na granularidade vendedor-pedido,
8,97\% das participações possuem ao menos um item fora do limite. A taxa
de atraso associado é 8,02\% e a nota média associada é 4,14.

\subsection{Validação e
controles}\label{docs__analytics__catalogo_metricas_analiticas.md__validauxe7uxe3o-e-controles}

{\def\LTcaptype{none} % do not increment counter
\begin{longtable}[]{@{}
  >{\raggedright\arraybackslash}p{(\linewidth - 4\tabcolsep) * \real{0.3333}}
  >{\raggedright\arraybackslash}p{(\linewidth - 4\tabcolsep) * \real{0.3333}}
  >{\raggedright\arraybackslash}p{(\linewidth - 4\tabcolsep) * \real{0.3333}}@{}}
\toprule\noalign{}
\begin{minipage}[b]{\linewidth}\raggedright
Risco
\end{minipage} & \begin{minipage}[b]{\linewidth}\raggedright
Controle
\end{minipage} & \begin{minipage}[b]{\linewidth}\raggedright
Resultado
\end{minipage} \\
\midrule\noalign{}
\endhead
\bottomrule\noalign{}
\endlastfoot
Multiplicação de itens e pagamentos &
\nolinkurl{queries/validation/02_controle_join_itens_pagamentos.sql} & O
join direto multiplica linhas e não pode sustentar somas. \\
Diferenças entre bruto e pagamento &
\nolinkurl{queries/validation/01_reconciliacao_financeira.sql} &
Reconciliação preservada por pedido e classificada com tolerância de R\$
0,01. \\
Cobertura de itens e pagamentos &
\nolinkurl{queries/validation/03_cobertura_relacoes_pedido.sql} &
Exceções permanecem visíveis e determinam a população de cada
métrica. \\
Atribuição a vendedores &
\nolinkurl{queries/validation/04_granularidade_vendedor.sql} & Itens e
valor atribuído reconciliam sem diferença; participações adicionais
refletem pedidos multivendedor. \\
Requisitos e relações &
\nolinkurl{queries/validation/05_cobertura_requisitos_funcionais.sql} &
Cobertura e limitações das relações estão documentadas por requisito. \\
\end{longtable}
}

\subsection{Decisões para ANALYTICS e
BI}\label{docs__analytics__catalogo_metricas_analiticas.md__decisuxf5es-para-analytics-e-bi}

\begin{itemize}
\tightlist
\item
  A base financeira por pedido deve estabilizar as pré-agregações de
  itens e pagamentos e informar flags de cobertura.
\item
  A estrutura de vendedor deve preservar as granularidades vendedor e
  vendedor-pedido, sem repetir o pagamento integral.
\item
  Métricas mensais devem expor a completude da cobertura antes de
  permitir comparação temporal.
\item
  Valores de pedidos cancelados devem permanecer em colunas separadas
  por representação financeira.
\item
  Nomes de campos de BI devem incluir a medida-base e, para vendedor, o
  termo \nolinkurl{associado} quando a informação pertencer ao pedido.
\item
  HHI e classe ABC são diagnósticos complementares à VEN04, não métricas
  de mercado nem critérios regulatórios.
\end{itemize}

\section{Estruturas do schema
ANALYTICS}\label{docs__analytics__estruturas_schema_analytics.md__estruturas-do-schema-analytics}

\subsection{Finalidade}\label{docs__analytics__estruturas_schema_analytics.md__finalidade}

As views do schema \nolinkurl{analytics} estabilizam combinações e
regras recorrentes validadas nas etapas M06-01 a M06-04. Elas reduzem
repetição, tornam explícita a granularidade de consumo e preservam
rastreabilidade até a CORE.

Não são cópias das tabelas de origem, não alteram a CORE e não
materializam dados. A seleção final de campos, relacionamentos e visuais
para Power BI pertence à etapa M06-06.

\subsection{Mapa de
dependências}\label{docs__analytics__estruturas_schema_analytics.md__mapa-de-dependuxeancias}

\begin{Shaded}
\begin{Highlighting}[]
\NormalTok{core.pedido ─┬─ core.cliente}
\NormalTok{             ├─ core.item\_pedido}
\NormalTok{             ├─ core.pagamento}
\NormalTok{             └─ core.avaliacao}
\NormalTok{                      │}
\NormalTok{                      ▼}
\NormalTok{        analytics.vw\_pedido\_financeiro}
\NormalTok{                      │}
\NormalTok{                      ▼}
\NormalTok{        analytics.vw\_financeiro\_mensal}

\NormalTok{core.item\_pedido ─┬─ core.pedido ─ core.cliente}
\NormalTok{                  ├─ core.vendedor}
\NormalTok{                  └─ core.avaliacao}
\NormalTok{                           │}
\NormalTok{                           ▼}
\NormalTok{           analytics.vw\_vendedor\_pedido}
\NormalTok{                           │}
\NormalTok{              core.produto ┤}
\NormalTok{                           ▼}
\NormalTok{        analytics.vw\_desempenho\_vendedor}
\end{Highlighting}
\end{Shaded}

\subsection{\texorpdfstring{\nolinkurl{analytics.vw_pedido_financeiro}}{}}\label{docs__analytics__estruturas_schema_analytics.md__analytics.vw_pedido_financeiro}

\textbf{Finalidade:} fornecer a base financeira e operacional
reutilizável de pedido, evitando o join direto entre itens, pagamentos e
avaliações.

\textbf{Granularidade:} uma linha por \nolinkurl{id_pedido}, incluindo
todos os status.

\textbf{População observada:} 99.441 linhas, sem duplicação de pedido.

{\def\LTcaptype{none} % do not increment counter
\begin{longtable}[]{@{}
  >{\raggedright\arraybackslash}p{(\linewidth - 4\tabcolsep) * \real{0.3333}}
  >{\raggedright\arraybackslash}p{(\linewidth - 4\tabcolsep) * \real{0.3333}}
  >{\raggedright\arraybackslash}p{(\linewidth - 4\tabcolsep) * \real{0.3333}}@{}}
\toprule\noalign{}
\begin{minipage}[b]{\linewidth}\raggedright
Grupo
\end{minipage} & \begin{minipage}[b]{\linewidth}\raggedright
Campos
\end{minipage} & \begin{minipage}[b]{\linewidth}\raggedright
Definição
\end{minipage} \\
\midrule\noalign{}
\endhead
\bottomrule\noalign{}
\endlastfoot
Identidade & \nolinkurl{id_pedido}, \nolinkurl{id_cliente},
\nolinkurl{id_cliente_unico} & Chaves do pedido e identidade persistente
do cliente. \\
Cliente & \nolinkurl{cidade_cliente}, \nolinkurl{estado_cliente} &
Localização cadastral, não endereço exato da entrega. \\
Pedido & \nolinkurl{status_pedido}, \nolinkurl{data_compra},
\nolinkurl{mes_compra}, \nolinkurl{data_aprovacao},
\nolinkurl{data_envio_transportador}, \nolinkurl{data_entrega},
\nolinkurl{data_estimada} & Estado e marcos temporais preservados da
CORE. \\
Entrega & \nolinkurl{entrega_atrasada} &
\nolinkurl{data_entrega > data_estimada}; nulo quando não calculável. \\
Cobertura & \nolinkurl{possui_itens}, \nolinkurl{possui_pagamentos},
\nolinkurl{possui_avaliacao} & Flags que impedem interpretar ausência
como valor zero. \\
Itens & \nolinkurl{quantidade_itens}, \nolinkurl{quantidade_vendedores},
\nolinkurl{valor_itens}, \nolinkurl{valor_frete},
\nolinkurl{valor_bruto} & Itens pré-agregados por pedido; valores ficam
nulos quando não há item. \\
Pagamentos & \nolinkurl{quantidade_pagamentos},
\nolinkurl{quantidade_tipos_pagamento}, \nolinkurl{maximo_parcelas},
\nolinkurl{valor_pago_registrado} & Sequências pré-agregadas por pedido;
valores ficam nulos quando não há pagamento. \\
Reconciliação & \nolinkurl{diferenca_reconciliacao} & Pagamento menos
bruto, somente quando ambas as relações existem. \\
Avaliação & \nolinkurl{quantidade_avaliacoes},
\nolinkurl{nota_media_pedido} & Avaliações pré-agregadas por pedido. \\
\end{longtable}
}

\subsection{\texorpdfstring{\nolinkurl{analytics.vw_financeiro_mensal}}{}}\label{docs__analytics__estruturas_schema_analytics.md__analytics.vw_financeiro_mensal}

\textbf{Finalidade:} disponibilizar os KPIs financeiros recorrentes por
mês de compra, incluindo cobertura temporal e cancelamentos.

\textbf{Granularidade:} uma linha por \nolinkurl{mes_compra}.

\textbf{População observada:} 25 meses.

{\def\LTcaptype{none} % do not increment counter
\begin{longtable}[]{@{}
  >{\raggedright\arraybackslash}p{(\linewidth - 4\tabcolsep) * \real{0.3333}}
  >{\raggedright\arraybackslash}p{(\linewidth - 4\tabcolsep) * \real{0.3333}}
  >{\raggedright\arraybackslash}p{(\linewidth - 4\tabcolsep) * \real{0.3333}}@{}}
\toprule\noalign{}
\begin{minipage}[b]{\linewidth}\raggedright
Grupo
\end{minipage} & \begin{minipage}[b]{\linewidth}\raggedright
Campos
\end{minipage} & \begin{minipage}[b]{\linewidth}\raggedright
Definição
\end{minipage} \\
\midrule\noalign{}
\endhead
\bottomrule\noalign{}
\endlastfoot
Tempo & \nolinkurl{mes_compra}, \nolinkurl{cobertura_diaria_completa} &
Mês de compra e sinalização de presença de pedidos em todos os dias do
mês. \\
Pedidos & \nolinkurl{pedidos}, \nolinkurl{pedidos_entregues},
\nolinkurl{pedidos_entregues_completos}, \nolinkurl{pedidos_cancelados}
& Volumes total, entregue, entregue com itens e pagamentos, e
cancelado. \\
Cancelamento & \nolinkurl{taxa_cancelamento_quantidade_percentual} &
Cancelados divididos por todos os pedidos do mês. \\
Financeiro entregue & \nolinkurl{valor_itens_entregues},
\nolinkurl{valor_frete_entregue}, \nolinkurl{valor_bruto_entregue},
\nolinkurl{valor_pago_registrado_entregue} & Somas sobre pedidos
entregues com itens e pagamentos. \\
Reconciliação & \nolinkurl{diferenca_reconciliacao_entregue} & Soma das
diferenças por pedido; não substitui o controle individual. \\
Médias e composição & \nolinkurl{valor_bruto_medio_pedido_entregue},
\nolinkurl{participacao_frete_percentual} & Média explícita do bruto e
participação do frete. \\
Evolução & \nolinkurl{crescimento_mensal_valor_itens_percentual} &
Calculado apenas quando mês atual e anterior têm cobertura diária
completa. \\
Cancelados & \nolinkurl{valor_itens_associado_cancelados},
\nolinkurl{valor_frete_associado_cancelados},
\nolinkurl{valor_bruto_associado_cancelados},
\nolinkurl{valor_pago_registrado_associado_cancelados} & Representações
separadas; não significam perda ou estorno. \\
\end{longtable}
}

\subsection{\texorpdfstring{\nolinkurl{analytics.vw_vendedor_pedido}}{}}\label{docs__analytics__estruturas_schema_analytics.md__analytics.vw_vendedor_pedido}

\textbf{Finalidade:} estabilizar a combinação entre contribuição direta
do vendedor e informações apenas associadas ao pedido.

\textbf{Granularidade:} uma linha por par \nolinkurl{id_vendedor},
\nolinkurl{id_pedido}, somente para pedidos entregues.

\textbf{População observada:} 97.819 linhas.

{\def\LTcaptype{none} % do not increment counter
\begin{longtable}[]{@{}
  >{\raggedright\arraybackslash}p{(\linewidth - 4\tabcolsep) * \real{0.3333}}
  >{\raggedright\arraybackslash}p{(\linewidth - 4\tabcolsep) * \real{0.3333}}
  >{\raggedright\arraybackslash}p{(\linewidth - 4\tabcolsep) * \real{0.3333}}@{}}
\toprule\noalign{}
\begin{minipage}[b]{\linewidth}\raggedright
Grupo
\end{minipage} & \begin{minipage}[b]{\linewidth}\raggedright
Campos
\end{minipage} & \begin{minipage}[b]{\linewidth}\raggedright
Definição
\end{minipage} \\
\midrule\noalign{}
\endhead
\bottomrule\noalign{}
\endlastfoot
Identidade e tempo & \nolinkurl{id_vendedor}, \nolinkurl{id_pedido},
\nolinkurl{data_compra}, \nolinkurl{mes_compra} & Participação do
vendedor no pedido e referência temporal. \\
Geografia & \nolinkurl{cidade_vendedor}, \nolinkurl{estado_vendedor},
\nolinkurl{cidade_cliente}, \nolinkurl{estado_cliente} & Origem
cadastral do vendedor e destino cadastral do cliente. \\
Itens diretos & \nolinkurl{quantidade_itens}, \nolinkurl{valor_itens},
\nolinkurl{valor_frete}, \nolinkurl{valor_bruto} & Medidas diretamente
atribuíveis aos itens do vendedor. \\
Envio direto & \nolinkurl{itens_com_envio_observado},
\nolinkurl{itens_enviados_apos_limite},
\nolinkurl{possui_item_enviado_apos_limite} & Contagens por item e flag
de presença no par vendedor-pedido. \\
Pedido associado & \nolinkurl{entrega_atrasada_associada},
\nolinkurl{nota_media_associada} & Entrega final e avaliação pertencem
ao pedido; não atribuem responsabilidade. \\
\end{longtable}
}

\subsection{\texorpdfstring{\nolinkurl{analytics.vw_desempenho_vendedor}}{}}\label{docs__analytics__estruturas_schema_analytics.md__analytics.vw_desempenho_vendedor}

\textbf{Finalidade:} fornecer uma visão histórica consolidada de
contribuição, concentração, alcance e qualidade associada dos
vendedores.

\textbf{Granularidade:} uma linha por vendedor com item em pedido
entregue.

\textbf{População observada:} 2.970 linhas.

{\def\LTcaptype{none} % do not increment counter
\begin{longtable}[]{@{}
  >{\raggedright\arraybackslash}p{(\linewidth - 4\tabcolsep) * \real{0.3333}}
  >{\raggedright\arraybackslash}p{(\linewidth - 4\tabcolsep) * \real{0.3333}}
  >{\raggedright\arraybackslash}p{(\linewidth - 4\tabcolsep) * \real{0.3333}}@{}}
\toprule\noalign{}
\begin{minipage}[b]{\linewidth}\raggedright
Grupo
\end{minipage} & \begin{minipage}[b]{\linewidth}\raggedright
Campos
\end{minipage} & \begin{minipage}[b]{\linewidth}\raggedright
Definição
\end{minipage} \\
\midrule\noalign{}
\endhead
\bottomrule\noalign{}
\endlastfoot
Identidade & \nolinkurl{id_vendedor}, \nolinkurl{cidade_vendedor},
\nolinkurl{estado_vendedor} & Vendedor e localização cadastral de
origem. \\
Ranking & \nolinkurl{posicao_valor_itens},
\nolinkurl{participacao_valor_itens_percentual},
\nolinkurl{participacao_acumulada_percentual}, \nolinkurl{classe_abc},
\nolinkurl{hhi_base_observada} & Concentração interna da base; não
representa mercado ou avaliação regulatória. \\
Volumes e valores & \nolinkurl{pedidos_com_participacao},
\nolinkurl{quantidade_itens}, \nolinkurl{valor_itens},
\nolinkurl{valor_frete}, \nolinkurl{valor_bruto} & Contribuição direta
histórica em pedidos entregues. \\
Médias & \nolinkurl{valor_medio_por_item},
\nolinkurl{valor_medio_por_pedido_com_participacao} & Médias com
numerador e denominador explícitos. \\
Diversidade e alcance & \nolinkurl{categorias_comercializadas},
\nolinkurl{estados_clientes_atendidos},
\nolinkurl{cidades_clientes_atendidas} & Diversidade e alcance
observados, não cobertura comercial formal. \\
Envio & \nolinkurl{itens_com_envio_observado},
\nolinkurl{itens_enviados_apos_limite},
\nolinkurl{taxa_itens_enviados_apos_limite_percentual} & VEN08 na
granularidade correta de item. \\
Entrega associada & \nolinkurl{pedidos_com_prazo_entrega_observado},
\nolinkurl{pedidos_atrasados_associados},
\nolinkurl{taxa_pedidos_atrasados_associados_percentual} & Atraso final
associado à participação do vendedor. \\
Avaliação associada & \nolinkurl{pedidos_com_avaliacao_associada},
\nolinkurl{nota_media_associada},
\nolinkurl{pedidos_com_avaliacao_negativa_associada},
\nolinkurl{taxa_avaliacao_negativa_associada_percentual} & Avaliações do
pedido, sem atribuição causal ao vendedor. \\
\end{longtable}
}

\subsection{Instalação e
validação}\label{docs__analytics__estruturas_schema_analytics.md__instalauxe7uxe3o-e-validauxe7uxe3o}

O \nolinkurl{database.setup} executa as views na ordem numérica de
dependência e valida a existência das quatro estruturas com
\nolinkurl{models/analytics/validate_views.sql}.

Após a carga, \nolinkurl{queries/validation/06_views_analytics.sql}
compara as granularidades e os valores atribuídos aos vendedores contra
a CORE. Na carga aprovada, todas as diferenças de contagem e valor
resultaram em zero.

As views são normais, não materializadas. Essa decisão evita duplicar
dados e criar política de atualização antes de existir evidência de
necessidade de desempenho. Materialização ou índices específicos
exigirão medição própria.

\subsection{Decisões e
limitações}\label{docs__analytics__estruturas_schema_analytics.md__decisuxf5es-e-limitauxe7uxf5es}

\begin{itemize}
\tightlist
\item
  Valores nulos em relações ausentes são preservados; contagens usam
  zero.
\item
  Pagamentos não são alocados aos vendedores.
\item
  A visão mensal principal compara itens e pagamentos somente na
  população entregue que possui ambas as relações.
\item
  O crescimento mensal permanece nulo em comparações com cobertura
  parcial.
\item
  A taxa VEN08 é calculada por item (9,32\%). O indicador exploratório
  anterior de 8,97\% representa pares vendedor-pedido com algum item
  fora do limite.
\item
  As estruturas são adequadas ao consumo SQL recorrente. O contrato
  final de importação, relacionamentos e medidas do Power BI está
  documentado em
  \href{consumo_power_bi_sql.md}{\nolinkurl{consumo_power_bi_sql.md}}.
\end{itemize}

\section{Contrato de consumo no Power BI e
SQL}\label{docs__analytics__consumo_power_bi_sql.md__contrato-de-consumo-no-power-bi-e-sql}

\subsection{Objetivo}\label{docs__analytics__consumo_power_bi_sql.md__objetivo}

Este documento consolida o contrato de consumo da M06-06. Ele organiza
as views aprovadas do schema \nolinkurl{analytics} como datasets para
exploração SQL e construção futura de um modelo semântico no Power BI,
sem copiar dados da CORE nem criar novas estruturas sem necessidade
demonstrada.

O contrato cobre nomes, granularidades, chaves, relacionamentos,
medidas, filtros, atualização, controles e limitações. A criação de um
arquivo \nolinkurl{.pbix}, de páginas ou de visuais não faz parte desta
entrega.

\subsection{Decisão de
arquitetura}\label{docs__analytics__consumo_power_bi_sql.md__decisuxe3o-de-arquitetura}

As quatro views criadas na M06-05 já estabilizam as regras recorrentes e
são as fontes canônicas de consumo. A M06-06 não cria marts adicionais
porque:

\begin{itemize}
\tightlist
\item
  as granularidades de pedido, mês, vendedor-pedido e vendedor já estão
  explícitas;
\item
  as combinações de itens, pagamentos e avaliações já são protegidas
  contra multiplicação de medidas;
\item
  não existe evidência de requisito de desempenho que justifique
  materialização;
\item
  novas cópias aumentariam o risco de divergência sem acrescentar
  semântica.
\end{itemize}

Para exportações controladas, \nolinkurl{queries/consumption/} expõe
listas explícitas de colunas. Consumidores com acesso direto ao
PostgreSQL podem consultar as views ou usar essas consultas como
instruções SQL no conector.

\subsection{Datasets}\label{docs__analytics__consumo_power_bi_sql.md__datasets}

{\def\LTcaptype{none} % do not increment counter
\begin{longtable}[]{@{}
  >{\raggedright\arraybackslash}p{(\linewidth - 8\tabcolsep) * \real{0.2000}}
  >{\raggedright\arraybackslash}p{(\linewidth - 8\tabcolsep) * \real{0.2000}}
  >{\raggedright\arraybackslash}p{(\linewidth - 8\tabcolsep) * \real{0.2000}}
  >{\raggedright\arraybackslash}p{(\linewidth - 8\tabcolsep) * \real{0.2000}}
  >{\raggedright\arraybackslash}p{(\linewidth - 8\tabcolsep) * \real{0.2000}}@{}}
\toprule\noalign{}
\begin{minipage}[b]{\linewidth}\raggedright
Nome recomendado no BI
\end{minipage} & \begin{minipage}[b]{\linewidth}\raggedright
Fonte canônica
\end{minipage} & \begin{minipage}[b]{\linewidth}\raggedright
Granularidade
\end{minipage} & \begin{minipage}[b]{\linewidth}\raggedright
Chave
\end{minipage} & \begin{minipage}[b]{\linewidth}\raggedright
Uso principal
\end{minipage} \\
\midrule\noalign{}
\endhead
\bottomrule\noalign{}
\endlastfoot
\nolinkurl{fato_pedido_financeiro} &
\nolinkurl{analytics.vw_pedido_financeiro} & Um pedido, incluindo todos
os status & \nolinkurl{id_pedido} & Detalhe financeiro, operacional e
reconciliação \\
\nolinkurl{agregado_financeiro_mensal} &
\nolinkurl{analytics.vw_financeiro_mensal} & Um mês de compra &
\nolinkurl{mes_compra} & Série executiva pronta para SQL e controles
mensais \\
\nolinkurl{fato_vendedor_pedido} &
\nolinkurl{analytics.vw_vendedor_pedido} & Um vendedor em um pedido
entregue & \nolinkurl{id_vendedor}, \nolinkurl{id_pedido} & Evolução,
alcance e risco associado ao vendedor \\
\nolinkurl{resumo_desempenho_vendedor} &
\nolinkurl{analytics.vw_desempenho_vendedor} & Um vendedor ativo no
histórico entregue & \nolinkurl{id_vendedor} & Ranking, curva ABC e
perfil histórico \\
\end{longtable}
}

Os valores de \nolinkurl{fato_vendedor_pedido} são contribuições diretas
dos itens do vendedor. Entrega e avaliação são apenas associadas ao
pedido. Pagamentos não são alocados ao vendedor.

\subsection{Perfil recomendado de
importação}\label{docs__analytics__consumo_power_bi_sql.md__perfil-recomendado-de-importauxe7uxe3o}

\subsubsection{Modelo semântico
principal}\label{docs__analytics__consumo_power_bi_sql.md__modelo-semuxe2ntico-principal}

Importar:

\begin{itemize}
\tightlist
\item
  \nolinkurl{fato_pedido_financeiro} para análises financeiras e
  operacionais;
\item
  \nolinkurl{fato_vendedor_pedido} para análises temporais de
  vendedores;
\item
  \nolinkurl{resumo_desempenho_vendedor} para ranking e segmentação
  histórica.
\end{itemize}

Criar no Power BI uma dimensão calendário contínua entre as datas mínima
e máxima de compra. A coluna de data deve ser relacionada, com direção
única, a \nolinkurl{data_compra} das duas tabelas fato. Para uso mensal,
a dimensão deve expor ano, mês, número do mês e uma chave
\nolinkurl{AAAA-MM} com ordenação cronológica.

Relacionamento permitido:

\begin{Shaded}
\begin{Highlighting}[]
\NormalTok{dim\_calendario[data]}
\NormalTok{    1 ─── * fato\_pedido\_financeiro[data\_compra\_data]}
\NormalTok{    1 ─── * fato\_vendedor\_pedido[data\_compra\_data]}

\NormalTok{resumo\_desempenho\_vendedor[id\_vendedor]}
\NormalTok{    1 ─── * fato\_vendedor\_pedido[id\_vendedor]}
\end{Highlighting}
\end{Shaded}

\nolinkurl{data_compra_data} representa uma coluna derivada somente com
a data de \nolinkurl{data_compra}. O filtro deve fluir da dimensão para
os fatos. Não usar direção bidirecional como configuração padrão.

\subsubsection{Agregado
mensal}\label{docs__analytics__consumo_power_bi_sql.md__agregado-mensal}

\nolinkurl{agregado_financeiro_mensal} é opcional no modelo Power BI.
Ele é recomendado para consulta SQL direta, validação e protótipos
executivos. Quando importado no mesmo modelo, deve permanecer separado
ou ser relacionado por uma dimensão de mês com chave única.

Não somar suas medidas com medidas calculadas a partir de
\nolinkurl{fato_pedido_financeiro}; ambas representam o mesmo domínio em
granularidades diferentes.

\subsubsection{Relacionamentos
proibidos}\label{docs__analytics__consumo_power_bi_sql.md__relacionamentos-proibidos}

\begin{itemize}
\tightlist
\item
  não relacionar diretamente \nolinkurl{fato_pedido_financeiro} e
  \nolinkurl{fato_vendedor_pedido} por \nolinkurl{id_pedido} para
  calcular valores financeiros;
\item
  não usar \nolinkurl{id_cliente} como dimensão persistente de pessoa; a
  identidade recorrente disponível é \nolinkurl{id_cliente_unico};
\item
  não combinar as tabelas por cidade sem estado;
\item
  não expandir geolocalização por prefixo de CEP, pois ela possui
  múltiplas ocorrências e multiplicaria linhas.
\end{itemize}

\subsection{Propriedade das
medidas}\label{docs__analytics__consumo_power_bi_sql.md__propriedade-das-medidas}

{\def\LTcaptype{none} % do not increment counter
\begin{longtable}[]{@{}
  >{\raggedright\arraybackslash}p{(\linewidth - 4\tabcolsep) * \real{0.3333}}
  >{\raggedright\arraybackslash}p{(\linewidth - 4\tabcolsep) * \real{0.3333}}
  >{\raggedright\arraybackslash}p{(\linewidth - 4\tabcolsep) * \real{0.3333}}@{}}
\toprule\noalign{}
\begin{minipage}[b]{\linewidth}\raggedright
Tema
\end{minipage} & \begin{minipage}[b]{\linewidth}\raggedright
Dataset proprietário
\end{minipage} & \begin{minipage}[b]{\linewidth}\raggedright
Regra
\end{minipage} \\
\midrule\noalign{}
\endhead
\bottomrule\noalign{}
\endlastfoot
Valores e contagens de pedido & \nolinkurl{fato_pedido_financeiro} &
Agregar uma vez por \nolinkurl{id_pedido} \\
Evolução financeira pronta & \nolinkurl{agregado_financeiro_mensal} &
Usar somente como agregado mensal \\
Valores e volumes do vendedor & \nolinkurl{fato_vendedor_pedido} & Somar
contribuições diretas do vendedor \\
Ranking e concentração histórica &
\nolinkurl{resumo_desempenho_vendedor} & Não recalcular como série
temporal \\
Pagamentos & \nolinkurl{fato_pedido_financeiro} & Nunca atribuir
integralmente ao vendedor \\
Entrega e avaliação por vendedor & \nolinkurl{fato_vendedor_pedido} &
Nomear como informação associada \\
\end{longtable}
}

Medidas de pedido devem ser calculadas em
\nolinkurl{fato_pedido_financeiro}; medidas de vendedor devem ser
calculadas em \nolinkurl{fato_vendedor_pedido}. Essa separação evita que
um pedido multivendedor repita valor bruto, pagamento ou avaliação.

\subsection{Medidas mínimas no Power
BI}\label{docs__analytics__consumo_power_bi_sql.md__medidas-muxednimas-no-power-bi}

As medidas abaixo são referências para o modelo semântico. Os
percentuais são razões e devem receber formatação percentual no Power
BI, sem multiplicação adicional por 100.

\begin{Shaded}
\begin{Highlighting}[]
\NormalTok{Pedidos =}
\NormalTok{DISTINCTCOUNT ( fato\_pedido\_financeiro[id\_pedido] )}

\NormalTok{Pedidos Entregues =}
\NormalTok{CALCULATE (}
\NormalTok{    [Pedidos],}
\NormalTok{    fato\_pedido\_financeiro[status\_pedido] = "delivered"}
\NormalTok{)}

\NormalTok{Pedidos Entregues Completos =}
\NormalTok{CALCULATE (}
\NormalTok{    [Pedidos Entregues],}
\NormalTok{    fato\_pedido\_financeiro[possui\_itens] = TRUE (),}
\NormalTok{    fato\_pedido\_financeiro[possui\_pagamentos] = TRUE ()}
\NormalTok{)}

\NormalTok{Valor dos Itens Entregues =}
\NormalTok{CALCULATE (}
\NormalTok{    SUM ( fato\_pedido\_financeiro[valor\_itens] ),}
\NormalTok{    fato\_pedido\_financeiro[status\_pedido] = "delivered",}
\NormalTok{    fato\_pedido\_financeiro[possui\_itens] = TRUE (),}
\NormalTok{    fato\_pedido\_financeiro[possui\_pagamentos] = TRUE ()}
\NormalTok{)}

\NormalTok{Valor Bruto Entregue =}
\NormalTok{CALCULATE (}
\NormalTok{    SUM ( fato\_pedido\_financeiro[valor\_bruto] ),}
\NormalTok{    fato\_pedido\_financeiro[status\_pedido] = "delivered",}
\NormalTok{    fato\_pedido\_financeiro[possui\_itens] = TRUE (),}
\NormalTok{    fato\_pedido\_financeiro[possui\_pagamentos] = TRUE ()}
\NormalTok{)}

\NormalTok{Valor Pago Registrado Entregue =}
\NormalTok{CALCULATE (}
\NormalTok{    SUM ( fato\_pedido\_financeiro[valor\_pago\_registrado] ),}
\NormalTok{    fato\_pedido\_financeiro[status\_pedido] = "delivered",}
\NormalTok{    fato\_pedido\_financeiro[possui\_itens] = TRUE (),}
\NormalTok{    fato\_pedido\_financeiro[possui\_pagamentos] = TRUE ()}
\NormalTok{)}

\NormalTok{Diferença de Reconciliação Entregue =}
\NormalTok{[Valor Pago Registrado Entregue] {-} [Valor Bruto Entregue]}

\NormalTok{Valor Bruto Médio por Pedido Entregue =}
\NormalTok{DIVIDE (}
\NormalTok{    [Valor Bruto Entregue],}
\NormalTok{    [Pedidos Entregues Completos]}
\NormalTok{)}

\NormalTok{Participação do Frete =}
\NormalTok{DIVIDE (}
\NormalTok{    CALCULATE (}
\NormalTok{        SUM ( fato\_pedido\_financeiro[valor\_frete] ),}
\NormalTok{        fato\_pedido\_financeiro[status\_pedido] = "delivered",}
\NormalTok{        fato\_pedido\_financeiro[possui\_itens] = TRUE (),}
\NormalTok{        fato\_pedido\_financeiro[possui\_pagamentos] = TRUE ()}
\NormalTok{    ),}
\NormalTok{    [Valor Bruto Entregue]}
\NormalTok{)}

\NormalTok{Taxa de Cancelamento =}
\NormalTok{VAR PedidosCancelados =}
\NormalTok{    CALCULATE (}
\NormalTok{        [Pedidos],}
\NormalTok{        REMOVEFILTERS ( fato\_pedido\_financeiro[status\_pedido] ),}
\NormalTok{        fato\_pedido\_financeiro[status\_pedido] = "canceled"}
\NormalTok{    )}
\NormalTok{VAR TodosPedidos =}
\NormalTok{    CALCULATE (}
\NormalTok{        [Pedidos],}
\NormalTok{        REMOVEFILTERS ( fato\_pedido\_financeiro[status\_pedido] )}
\NormalTok{    )}
\NormalTok{RETURN}
\NormalTok{    DIVIDE ( PedidosCancelados, TodosPedidos )}

\NormalTok{Valor dos Itens por Vendedor =}
\NormalTok{SUM ( fato\_vendedor\_pedido[valor\_itens] )}

\NormalTok{Taxa de Itens Enviados Após o Limite =}
\NormalTok{DIVIDE (}
\NormalTok{    SUM ( fato\_vendedor\_pedido[itens\_enviados\_apos\_limite] ),}
\NormalTok{    SUM ( fato\_vendedor\_pedido[itens\_com\_envio\_observado] )}
\NormalTok{)}
\end{Highlighting}
\end{Shaded}

O crescimento mensal deve ser calculado apenas entre meses de cobertura
diária completa. A definição canônica já está em
\nolinkurl{agregado_financeiro_mensal[crescimento_mensal_valor_itens_percentual]};
reimplementações em DAX devem reproduzir essa condição.

\subsection{Tipos e
formatação}\label{docs__analytics__consumo_power_bi_sql.md__tipos-e-formatauxe7uxe3o}

{\def\LTcaptype{none} % do not increment counter
\begin{longtable}[]{@{}ll@{}}
\toprule\noalign{}
Grupo & Tratamento recomendado \\
\midrule\noalign{}
\endhead
\bottomrule\noalign{}
\endlastfoot
Identificadores & Texto; nunca resumir \\
Datas e timestamps & Data ou data/hora conforme a coluna original \\
Valores monetários & Número decimal fixo, moeda BRL \\
Contagens & Número inteiro \\
Flags & Booleano \\
Taxas das views & Número decimal já expresso em pontos percentuais \\
Taxas DAX com \nolinkurl{DIVIDE} & Número decimal formatado como
percentual \\
Notas & Número decimal, sem formatação percentual \\
\end{longtable}
}

As colunas percentuais das views SQL já estão multiplicadas por 100. Por
exemplo, \nolinkurl{9.32} significa \nolinkurl{9,32%}. Medidas DAX
baseadas em \nolinkurl{DIVIDE} retornam \nolinkurl{0.0932} e devem
apenas receber o formato percentual.

\subsection{Filtros e
interpretação}\label{docs__analytics__consumo_power_bi_sql.md__filtros-e-interpretauxe7uxe3o}

\begin{itemize}
\tightlist
\item
  a data principal é \nolinkurl{data_compra};
\item
  a visão financeira principal usa pedidos \nolinkurl{delivered} com
  itens e pagamentos;
\item
  meses sem cobertura diária completa não sustentam comparações de
  crescimento;
\item
  valores associados a cancelados não representam perda, estorno ou
  receita;
\item
  \nolinkurl{cidade_cliente} e \nolinkurl{estado_cliente} representam
  localização cadastral, não o endereço exato da entrega;
\item
  alcance geográfico é observado na base, não cobertura comercial
  formal;
\item
  atraso e avaliação associados ao vendedor não demonstram
  responsabilidade ou causalidade;
\item
  HHI e participação representam somente a base histórica observada.
\end{itemize}

\subsection{Atualização e
segurança}\label{docs__analytics__consumo_power_bi_sql.md__atualizauxe7uxe3o-e-seguranuxe7a}

As views são normais e refletem o estado corrente da CORE. A sequência
de atualização é:

\begin{Shaded}
\begin{Highlighting}[]
\NormalTok{CSV → RAW → CORE → validações do ELT → views ANALYTICS → atualização do BI}
\end{Highlighting}
\end{Shaded}

O Power BI deve usar uma credencial PostgreSQL de leitura com acesso
apenas aos objetos necessários. \nolinkurl{DATABASE_URL}, senhas e
arquivos locais não devem ser incorporados ao repositório ou às
consultas versionadas.

O modo \textbf{Importação} é o padrão recomendado para esta base
histórica. O uso de DirectQuery, atualização incremental, agregações
automáticas ou views materializadas depende de medição de desempenho e
de um requisito operacional futuro.

\subsection{Controles antes da
publicação}\label{docs__analytics__consumo_power_bi_sql.md__controles-antes-da-publicauxe7uxe3o}

Antes de atualizar ou publicar um modelo semântico:

\begin{enumerate}
\def\labelenumi{\arabic{enumi}.}
\tightlist
\item
  executar o ELT e a validação independente;
\item
  executar \nolinkurl{queries/validation/06_views_analytics.sql};
\item
  confirmar unicidade de \nolinkurl{id_pedido}, \nolinkurl{mes_compra} e
  \nolinkurl{id_vendedor} nas fontes correspondentes;
\item
  confirmar unicidade do par \nolinkurl{id_vendedor},
  \nolinkurl{id_pedido};
\item
  comparar totais financeiros do BI com
  \nolinkurl{vw_financeiro_mensal};
\item
  verificar tipos, formatos, direção de filtros e relacionamentos
  ativos;
\item
  registrar data da atualização e período coberto.
\end{enumerate}

\subsection{Limitações e
evolução}\label{docs__analytics__consumo_power_bi_sql.md__limitauxe7uxf5es-e-evoluuxe7uxe3o}

Este contrato prepara o consumo, mas não entrega monitoramento em tempo
real, governança de workspace, gateway, política de acesso, atualização
agendada ou design visual. Esses elementos dependem do ambiente Power BI
de destino.

Também não há evidência atual para materializar views ou criar índices
específicos. Qualquer evolução deve partir de medição reproduzível e
preservar a rastreabilidade até as tabelas CORE.

\end{document}
